\documentclass[11pt,reqno]{amsart}

\usepackage[T1]{fontenc}
\usepackage[utf8]{inputenc}
\usepackage{lmodern}
\usepackage{microtype}
\usepackage{mathtools}
\usepackage{amssymb}
\usepackage{mathrsfs}
\usepackage{enumitem}
\usepackage{aliascnt}
\usepackage{xcolor}
\usepackage[margin=1.12in]{geometry}
\usepackage[colorlinks=true,linkcolor=blue!45!black,citecolor=blue!45!black,
  urlcolor=blue!45!black,
  pdftitle={Theta--Heisenberg Untwisting, Uniform Prill Covers, and Abelian Brill--Noether Blocks},
  pdfauthor={OpenAI},
  pdfsubject={Algebraic geometry and finite Heisenberg representations},
  pdfkeywords={Raynaud bundle, finite Heisenberg group, theta group, Prill problem, Brill--Noether locus}]{hyperref}

\setlist{itemsep=2pt,topsep=4pt}
\numberwithin{equation}{section}

\newtheorem{theorem}{Theorem}[section]
\newaliascnt{proposition}{theorem}
\newtheorem{proposition}[proposition]{Proposition}
\aliascntresetthe{proposition}
\newaliascnt{lemma}{theorem}
\newtheorem{lemma}[lemma]{Lemma}
\aliascntresetthe{lemma}
\newaliascnt{corollary}{theorem}
\newtheorem{corollary}[corollary]{Corollary}
\aliascntresetthe{corollary}
\newaliascnt{definition}{theorem}
\newtheorem{definition}[definition]{Definition}
\aliascntresetthe{definition}
\newaliascnt{question}{theorem}
\newtheorem{question}[question]{Question}
\aliascntresetthe{question}
\theoremstyle{remark}
\newaliascnt{example}{theorem}
\newtheorem{example}[example]{Example}
\aliascntresetthe{example}
\newaliascnt{remark}{theorem}
\newtheorem{remark}[remark]{Remark}
\aliascntresetthe{remark}

\usepackage[nameinlink,noabbrev]{cleveref}
\crefname{theorem}{theorem}{theorems}
\crefname{proposition}{proposition}{propositions}
\crefname{lemma}{lemma}{lemmas}
\crefname{corollary}{corollary}{corollaries}
\crefname{definition}{definition}{definitions}
\crefname{question}{question}{questions}
\crefname{example}{example}{examples}
\crefname{remark}{remark}{remarks}

\DeclareMathOperator{\Pic}{Pic}
\DeclareMathOperator{\Nm}{Nm}
\DeclareMathOperator{\Ind}{Ind}
\DeclareMathOperator{\mult}{mult}
\DeclareMathOperator{\rk}{rk}
\DeclareMathOperator{\im}{im}
\DeclareMathOperator{\Hom}{Hom}
\DeclareMathOperator{\Ext}{Ext}
\DeclareMathOperator{\Spec}{Spec}
\DeclareMathOperator{\Cliff}{Cliff}
\DeclareMathOperator{\Gal}{Gal}
\DeclareMathOperator{\sign}{sign}
\newcommand{\cO}{\mathcal O}
\newcommand{\cV}{\mathcal V}
\newcommand{\cG}{\mathcal G}
\newcommand{\bC}{\mathbf C}
\newcommand{\bG}{\mathbf G}
\newcommand{\bZ}{\mathbf Z}
\newcommand{\bQ}{\mathbf Q}
\newcommand{\bP}{\mathbf P}
\newcommand{\whA}{\widehat A}

\title[Theta--Heisenberg Untwisting]{Theta--Heisenberg Untwisting, Uniform Prill Covers,\\
and Abelian Brill--Noether Blocks}
\author{OpenAI}
\date{}

\subjclass[2020]{14H30, 14H40, 14K25, 14D20, 20C15, 57K20}
\keywords{Raynaud bundle, finite Heisenberg group, theta group, Prill's
problem, Brill--Noether locus, Putman--Wieland conjecture}

\begin{document}

\begin{abstract}
For integers \(g\geq2\) and \(n\geq2\) with \(n\mid g\), I construct a
fixed finite unitary representation
\[
 \rho_{g,n}\colon\pi_1(\Sigma_g)\longrightarrow SU(n^g)
\]
such that, for every marked complex structure \(C\) on \(\Sigma_g\) and
every \(L\in\Pic^{g/n}(C)\), the associated flat bundle satisfies
\(H^0(C,\cV_{\rho_{g,n},C}\otimes L)\neq0\).  The construction identifies
the projective Narasimhan--Seshadri monodromy of a Raynaud bundle with the
Schr\"odinger representation of a finite theta group.

I then prove a theta--Heisenberg untwisting theorem.  The rank-\(n^g\)
bundle is realized as the direct image of an exact-order-\(n\) torsion
line bundle from a degree-\(n^g\) cover; a cyclic cover removes this
torsion and produces a
connected ordinary cover of degree \(n^{g+1}\); its Galois closure has
degree \(n^{2g+1}\).  These three degrees are sharp among covers retaining
the fixed Schr\"odinger block.  Taking \(n=g\) gives, on every curve of
genus \(g\), a connected \'etale cover of degree \(g^{g+1}\) for which
every fibre moves in a pencil.  In genus two the degree is \(8\).

To quantify the remaining gap toward global minimality, I prove a
strict rank--corank inequality: a nontrivial
stable degree-zero bundle \(W\), with
\[
 c=\min_p h^0(W(p))>0,
\]
satisfies \(\rk W\geq gc+2\).  This gives the lower bounds
\(\deg f\geq gk+3\) for covers with at least \(k+1\) fibre sections and
\(|\Gal(f)|\geq(g+2)(g+3)\) in the Galois Prill-exceptional case.
In genus two I determine the minimum over all genus-two curves.  On the curve
\(y^2=x^5-1\), a \(C_5\)-equivariant local system with natural
\(A_5\)-monodromy has \(w_2=1\), while an equivariant Lefschetz
calculation forces its two-dimensional theta obstruction to vanish.
It therefore produces a connected Prill-exceptional cover of degree
five; the strict rank bound shows that degree five is minimal.
The same method extends to Fourier--Mukai bundles pulled back from
arbitrary abelian quotients of a Jacobian, producing
positive-dimensional abelian blocks in classical Brill--Noether loci.
\end{abstract}

\maketitle

\section{Introduction}

Let \(C\) be a smooth projective complex curve and let
\(f\colon X\to C\) be a connected finite \'etale cover.  The fibre over
\(p\in C\) moves in a pencil precisely when
\[
 h^0\!\left(X,\cO_X(f^{-1}(p))\right)\geq2.
\]
Because this dimension is upper semicontinuous in \(p\), its generic
value is its minimum.  Thus the condition for a general fibre is
equivalent, in this setting, to the same inequality for every fibre.
Prill's problem asks whether this can occur for a general fibre
\cite[Chapter VI, Exercise D, p.~268]{ACGH}.  A
degree-\(36\) construction on every genus-two curve was obtained by
Landesman and Litt \cite{LandesmanLittPrill}.  The higher-genus question
was subsequently recorded together with conjectures on generic global
generation and horizontal generic vanishing in \cite[\S6.3]{LittMotives}.

I prove that every fibre can move, simultaneously for every marked
complex structure and in every genus \(g\geq2\).  The construction is
explicit and topological.  It does not arise from a single isolated
algebraic cover, but from a fixed finite local system whose cohomological
support locus is the whole Picard variety.

\begin{theorem}[Uniform Heisenberg--Raynaud local system]\label{thm:intro-uniform}
Let \(g\geq2\), let \(n\geq2\) divide \(g\), and put
\[
 m=\frac gn,\qquad r=n^g.
\]
There is a fixed irreducible representation
\[
 \rho_{g,n}\colon\pi_1(\Sigma_g)\longrightarrow SU(r)
\]
with image of order \(n^{2g+1}\) such that, for every marked genus-\(g\)
curve \(C\),
\[
 H^0(C,\cV_{\rho_{g,n},C}\otimes L)\neq0
 \quad\text{for every }L\in\Pic^m(C).
 \tag{1.1}\label{eq:intro-full-support}
\]
The bundle \(\cV_{\rho_{g,n},C}\) is stable of degree zero, has trivial
determinant, and has no global sections.
\end{theorem}

The divisibility \(n\mid g\) has two simultaneous meanings.  It makes the
Raynaud slope \(g/n\) integral, so that the Raynaud bundle can be
normalized to degree zero by a line bundle.  It also kills the scalar
obstruction to lifting the theta-group projective representation across
the surface relation.  Thus the numerical slope and the topological
surface obstruction are the same integrality condition.

Lagrangian induction then realizes the high-rank cohomological witness
as the direct image of a line bundle on a smaller cover; a cyclic torsor
removes its residual scalar monodromy.

\begin{theorem}[Theta--Heisenberg untwisting]\label{thm:intro-untwist}
Under the hypotheses of \cref{thm:intro-uniform}, there is a factorization
\[
 X\xrightarrow{\,v\,}Y\xrightarrow{\,u\,}C
\]
with the following properties.
\begin{enumerate}[label=\textup{(\roman*)}]
\item The cover \(u\) is connected Galois of degree \(r=n^g\).  There is
an exact-order-\(n\) torsion line bundle \(\lambda\in\Pic^0(Y)\) such that
\[
 u_*\lambda\simeq\cV_{\rho_{g,n},C},\qquad
 H^i(Y,\lambda\otimes u^*L)
 \simeq H^i(C,\cV_{\rho_{g,n},C}\otimes L)
 \tag{1.2}\label{eq:intro-exact-line}
\]
for every line bundle \(L\) and \(i=0,1\).
\item The map \(v\) is the connected cyclic degree-\(n\) cover
trivializing \(\lambda\).  Thus \(f=u\circ v\) has degree \(n^{g+1}\), and
\[
 h^0(X,f^*L)\geq h^0(C,L)+
 h^0(C,\cV_{\rho_{g,n},C}\otimes L).
 \tag{1.3}\label{eq:intro-section-bound}
\]
\item The Galois closure of \(f\) has degree \(n^{2g+1}\).
\end{enumerate}
The degrees \(r\), \(nr\), and \(nr^2\) are respectively minimal among
connected twisted-line, ordinary, and Galois covers which retain the
fixed Schr\"odinger local system.
\end{theorem}

For \(n=g\), equation \eqref{eq:intro-section-bound} applied to
\(L=\cO_C(p)\) gives the following extension of the genus-two result to
every genus.

\begin{corollary}[Uniform Prill cover]\label{cor:intro-prill}
For every \(g\geq2\) there is a fixed finite-index subgroup of
\(\pi_1(\Sigma_g)\) such that, for every marked complex structure
\(C\), the induced connected \'etale cover
\[
 f_{g,C}\colon X_{g,C}\longrightarrow C,
 \qquad \deg f_{g,C}=g^{g+1},
\]
satisfies
\[
 h^0\!\left(X_{g,C},\cO_{X_{g,C}}(f_{g,C}^{-1}(p))\right)\geq2
 \quad\text{for every }p\in C
\]
The topological cover is independent of \(C\).  It is non-Galois, and
its Galois closure has degree \(g^{2g+1}\).
\end{corollary}

The number \(g^{g+1}\) is not asserted to be the minimum among all
Prill-exceptional covers.  It is the exact minimum among transitive
covers whose permutation local system contains the particular
Raynaud--Schr\"odinger block in \cref{thm:intro-uniform}.  The distinction
is essential.  I prove the sharper multiplicity law
\[
 \deg f=n^{g+1}\,
 \mult_{\cV_{\rho_{g,n}}}(f_*\cO_X)
 \tag{1.4}\label{eq:intro-pareto}
\]
for covers obtained from transitive sets of the finite Heisenberg group
for which the displayed multiplicity is positive, and the corresponding
inequality for arbitrary connected covers retaining the block.

The construction has Brill--Noether consequences.  If
\(d=n^{g+1}\), then
\[
 f^*\Pic^m(C)\subset W_{dm}^0(X),\qquad
 f^*C^{(m)}\subset W_{dm}^1(X).
\]
For \(n\geq3\), the second locus has Brill--Noether number
\[
 \rho\bigl(g(X),1,dm\bigr)=d(2m-g+1)-3<0,
\]
although it contains an \(m\)-dimensional subvariety.  More generally,
intermediate isotropic quotients give an exact degree--multiplicity
trade-off and higher-rank Brill--Noether strata.

I also obtain a lower bound that does not use finite monodromy.

\begin{theorem}[Strict rank--corank inequality]\label{thm:intro-rank}
Let \(W\) be a nontrivial stable vector bundle of degree zero and rank
\(r\) on a curve of genus \(g\geq2\).  If
\[
 c=\min_{p\in C}h^0(C,W(p))\geq1,
\]
then
\[
 r\geq gc+2.
\tag{1.5}\label{eq:intro-rank}
\]
\end{theorem}

Consequently, if every fibre of a connected \'etale cover has at least
\(k+1\) sections, then its degree is at least \(gk+3\).  A Galois
Prill-exceptional cover has group order at least
\((g+2)(g+3)\).  Together with the implication that a Putman--Wieland
counterexample is Prill-exceptional
\cite[Lemma 5.5]{LandesmanLittApplications}, this excludes closed
Putman--Wieland counterexamples in these numerical ranges.  It does not
give the converse implication and does not settle the Putman--Wieland
conjecture outside those ranges.

A later section extends the method away from the full Jacobian.  An
ample line bundle on an abelian quotient of a Jacobian supplies a stable
semihomogeneous bundle, a Lagrangian line descent, and a finite
projective local system.  When its slope is integral, Heisenberg
untwisting produces a classical rank-one Brill--Noether block on a
finite cover.  Applied to a Prym subvariety, this gives a
Prym-direction support block; it does not by itself give a finite orbit
in Prym cohomology.

\section{Finite Heisenberg groups and scalar lifting}

\subsection{The standard Schr\"odinger representation}

Fix a primitive \(n\)-th root of unity \(\zeta\).  Define
\[
 H_{g,n}=
 \left\langle z,x_1,\ldots,x_g,y_1,\ldots,y_g\ \middle|\
 \begin{array}{c}
 z\text{ central},\quad z^n=x_i^n=y_i^n=1,\\
 {}[x_i,x_j]=[y_i,y_j]=1,\quad [x_i,y_j]=z^{\delta_{ij}}
 \end{array}
 \right\rangle,
\tag{2.1}\label{eq:heis-presentation}
\]
where \([a,b]=aba^{-1}b^{-1}\).  The relations collect every word into
the following normal form; uniqueness follows from the faithful
shift--clock action defined below.  Thus every element has a unique
expression
\[
 z^c x_1^{u_1}\cdots x_g^{u_g}
 y_1^{v_1}\cdots y_g^{v_g},
\]
so
\[
 |H_{g,n}|=n^{2g+1},\qquad Z(H_{g,n})=\langle z\rangle.
\tag{2.2}\label{eq:heis-order}
\]

Let \(T=(\bZ/n)^g\), let \(S=\bC[T]\), and write \(e_t\) for its standard
basis.  I use the central-weight convention dictated by the dual
theta representation in \cref{eq:schneider-splitting} and define
\[
 x_i e_t=e_{t+e_i},\qquad
 y_i e_t=\zeta^{t_i}e_t,\qquad
 z e_t=\zeta^{-1}e_t.
\tag{2.3}\label{eq:schrodinger}
\]
These matrices satisfy \eqref{eq:heis-presentation}.  The resulting
representation
\[
 \tau_{g,n}\colon H_{g,n}\longrightarrow U(S)
\]
is faithful and irreducible, and \(\dim S=r=n^g\).  I call \(S\) the
geometric Schr\"odinger module.  Its dual is the usual weight-one
module; fixing the convention here removes any later \(S/S^\vee\)
ambiguity.

\begin{lemma}[Character and determinant]\label{lem:character}
The character of \(S\) is
\[
 \chi_S(z^a)=r\zeta^{-a},\qquad
 \chi_S(h)=0\quad(h\notin\langle z\rangle).
\tag{2.4}\label{eq:sch-character}
\]
If \(g\geq2\), then \(\tau_{g,n}(H_{g,n})\subset SU(S)\).
\end{lemma}

\begin{proof}
If the shift part of \(h\) is nonzero, its matrix has no diagonal
entries.  If the shift part vanishes but the clock part is nonzero, its
trace is a product containing a complete nontrivial geometric sum.  This
proves \eqref{eq:sch-character}, while irreducibility follows at once
from the simultaneous clock eigenspaces and the transitive action of the
shifts.

Each \(x_i\) is a product of \(n^{g-1}\) cycles of length \(n\), and the
eigenvalues of \(y_i\) contain every \(n\)-th root with multiplicity
\(n^{g-1}\).  Hence
\[
 \det x_i=\det y_i=(-1)^{(n-1)n^{g-1}}=1.
\]
For even \(n\), the last equality uses \(g\geq2\); for odd \(n\) it is
immediate.  Finally, \(\det(\zeta^{-1}I)=\zeta^{-n^g}=1\).
\end{proof}

Set
\[
 A_{g,n}=H_{g,n}/\langle z\rangle\simeq(\bZ/n)^{2g}.
\tag{2.5}\label{eq:heisenberg-quotient}
\]
The scalar commutator in the geometric Schr\"odinger module induces a
perfect alternating pairing
\[
 e_n\colon A_{g,n}\times A_{g,n}\longrightarrow\mu_n,
 \qquad e_n(\bar x_i,\bar y_j)=\zeta^{-\delta_{ij}}.
\tag{2.6}\label{eq:weil-pairing}
\]
It is the inverse of the conventional Weil pairing attached to the dual
section module.  The finite Stone--von Neumann theorem says that the
fixed-central-weight irreducible projective representation is determined,
up to projective equivalence, by this pairing
\cite{MumfordTheta,Schneider}.

\subsection{Sharpness and the degree--multiplicity law}

For the remainder of this subsection, assume \(n\mid g\), put
\(m=g/n\), write
\[
 \pi_g=\left\langle a_1,b_1,\ldots,a_g,b_g\ \middle|\
 \prod_{i=1}^g[a_i,b_i]=1\right\rangle,
\]
and use the surface homomorphism
\[
 \varphi_{g,n}\colon\pi_g\longrightarrow H_{g,n},\qquad
 a_i\longmapsto x_i,\quad b_i\longmapsto y_i.
\]
The relation holds because its image is \(z^g=1\).  Let
\(\rho_{g,n}=\tau_{g,n}\circ\varphi_{g,n}\), and let
\(\cV_{\rho_{g,n},C}\) denote its flat bundle on a marked curve.  The
theta-theoretic support property of this bundle is proved in
\cref{thm:uniform}; the degree statements below are purely
representation-theoretic.

The next theorem makes the minimality statement precise without making a
claim about unrelated Prill constructions.

\begin{theorem}[Three sharp degree bounds]\label{thm:three-optima}
Let \(V=\cV_{\rho_{g,n},C}\), let \(r=n^g\), and let
\(\varphi=\varphi_{g,n}\).
\begin{enumerate}[label=\textup{(\roman*)}]
\item Suppose \(u'\colon Y'\to C\) is connected and \(\eta\) is a
rank-one finite local system on \(Y'\).  If \(V\) is a summand of
\(u'_*\eta\), then \(\deg u'\geq r\).
\item Suppose \(f'\colon X'\to C\) is connected and \(V\) occurs in
\(f'_*\cO_{X'}\) with multiplicity \(q>0\).  Then
\[
 \deg f'\geq nrq=n^{g+1}q.
\tag{2.7}\label{eq:arbitrary-cover-bound}
\]
If \(f'\) is the cover associated with a transitive \(H_{g,n}\)-set,
equality holds:
\[
 \deg f'=nrq.
\tag{2.8}\label{eq:exact-pareto}
\]
\item Suppose \(f'\) is connected Galois and \(V\) is a summand of
\(f'_*\cO_{X'}\).  Then \(\deg f'\geq nr^2=n^{2g+1}\).
\end{enumerate}
The covers in \cref{thm:exact-untwisting} attain the three respective
lower bounds.
\end{theorem}

\begin{proof}
Let \(\Gamma\subset\pi_g\) be the subgroup defining a connected cover and
put \(J=\varphi(\Gamma)\subset H=H_{g,n}\).  The map of coset sets
\(\pi_g/\Gamma\to H/J\) is surjective, so
\[
 [\pi_g:\Gamma]\geq[H:J].
\tag{2.9}\label{eq:index-comparison}
\]

For (i), Frobenius reciprocity gives a one-dimensional character of
\(\Gamma\) occurring in \(S|_\Gamma\).  Its restriction to
\(\Gamma\cap\ker\varphi\) is trivial because that subgroup acts as the
identity on \(S\); hence it descends to a character of \(J\).
Commutators in \(J\) act trivially on this line.  The faithful central
character of \(S\) therefore shows that the image of \(J\) in
\(A_{g,n}\) is isotropic.  A subgroup of a finite symplectic module of
order \(r^2\) has order at most \(r\).  Writing \(\bar J\) for that
image gives
\[
 |J|=|J\cap\langle z\rangle|\,|\bar J|\leq nr.
\]
Thus \([H:J]\geq r\).

For (ii), Frobenius reciprocity gives
\[
 q=\dim(S^\vee)^J.
\tag{2.10}\label{eq:multiplicity-invariants}
\]
This dimension is positive only if \(J\cap\langle z\rangle=1\).
Using \cref{lem:character},
\[
 q=\frac1{|J|}\sum_{j\in J}\chi_{S^\vee}(j)
 =\frac r{|J|}.
\tag{2.11}\label{eq:average-character}
\]
Consequently,
\[
 [H:J]=\frac{nr^2}{|J|}=nrq.
\]
Together with \eqref{eq:index-comparison}, this proves
\eqref{eq:arbitrary-cover-bound}; equality holds when
\(\Gamma=\varphi^{-1}(J)\).

For (iii), \(\Gamma\) is normal, so \(J\) is normal in \(H\).  The
nonzero invariant space \((S^\vee)^J\) is then \(H\)-stable.
Irreducibility gives \((S^\vee)^J=S^\vee\), and faithfulness gives
\(J=1\).  Equation
\eqref{eq:index-comparison} now yields \(\deg f'\geq|H|=nr^2\).
\end{proof}

\begin{remark}[Why the projective orbit is too small]\label{rem:projective-stabilizer}
The stabilizer \(B=\langle z\rangle K\) has index \(r\), but its ordinary
permutation module factors through \(H/\langle z\rangle\).  The centre
therefore acts trivially on every constituent, so that permutation
module cannot contain \(S\).  The degree-\(r\) stage retains \(S\) only
after inserting the character line \(\lambda\).  Removing that scalar
costs exactly the cyclic factor \(n\).
\end{remark}

\subsection{The isotropic tower}

For \(0\leq k\leq g\), set
\[
 K_k=\langle y_1,\ldots,y_k\rangle,\qquad
 \Gamma_k=\varphi_{g,n}^{-1}(K_k),
\]
and let \(f_k\colon X_k\to C\) be the associated cover.  Define
\[
 d_k=n^{2g+1-k},\qquad \mu_k=n^{g-k}.
\tag{2.12}\label{eq:tower-numbers}
\]

\begin{proposition}[Exact compression tower]\label{prop:compression-tower}
The covers \(f_k\) are connected and satisfy
\[
 \deg f_k=d_k,\qquad
 \mult_S\bC[H_{g,n}/K_k]=\mu_k,\qquad
 d_k=n^{g+1}\mu_k.
\tag{2.13}\label{eq:tower-pareto}
\]
There is a tower
\[
 X_0\longrightarrow X_1\longrightarrow\cdots
 \longrightarrow X_g\longrightarrow C
\tag{2.14}\label{eq:cover-tower}
\]
in which every map \(X_{k-1}\to X_k\) is cyclic \'etale of degree \(n\).
The cover \(X_0\to C\) is Galois with group \(H_{g,n}\); every
\(X_k\to C\) for \(k>0\) is non-Galois and has the same Galois closure.
Moreover,
\[
 (f_k)_*\cO_{X_k}
 \simeq\cO_C\oplus
 \cV_{\rho_{g,n},C}^{\oplus\mu_k}\oplus E_k
\tag{2.15}\label{eq:tower-splitting}
\]
and, for every \(L\in\Pic^m(C)\),
\[
 h^0(X_k,f_k^*L)\geq h^0(C,L)+\mu_k.
\tag{2.16}\label{eq:tower-sections}
\]
\end{proposition}

\begin{proof}
The order and multiplicity statements follow from
\eqref{eq:average-character}.  The inclusions
\(K_{k-1}\triangleleft K_k\) give the cyclic maps in
\eqref{eq:cover-tower}.  Conjugating a nontrivial element of \(K_k\) by
the corresponding \(x_i\) produces a nontrivial central factor; hence
\(K_k\) is not normal for \(k>0\), and the same argument as in
\cref{thm:exact-untwisting} shows that its core is trivial.
The representation splitting and \eqref{eq:tower-sections} follow from
Frobenius reciprocity, the projection formula, and
\cref{thm:uniform}.
\end{proof}

\begin{corollary}[Prill's problem in every genus]\label{cor:prill-all-genus}
Taking \(n=g\) and \(k=g\) gives a fixed topological cover whose
holomorphic realization
\(f_{g,C}\colon X_{g,C}\to C\), for every marked \(C\), is connected
and \'etale of degree
\[
 g^{g+1}
\]
such that every fibre moves in a pencil.  It is fixed by the marking and
works for every complex structure.  Its source has genus
\[
 1+g^{g+1}(g-1).
\tag{2.17}\label{eq:prill-source-genus}
\]
\end{corollary}

\begin{proof}
Here \(m=1\) and \(\mu_g=1\).  Apply
\eqref{eq:tower-sections} to \(L=\cO_C(p)\) and use Riemann--Hurwitz.
\end{proof}

\begin{example}[Genus two]\label{ex:genus-two}
For \(g=n=2\), put
\(K=K_g=\langle y_1,y_2\rangle\).  The \(K\)-quotient cover, which is minimal among ordinary
covers retaining the fixed Schr\"odinger block, has degree \(8\) and
source genus \(9\); its Galois closure has degree \(32\).  If
\(\alpha_1,\alpha_2,\alpha_3\) are the three nontrivial characters
appearing in the annihilator of \(K\), then
\[
 f_*\cO_X\simeq
 \cO_C\oplus\alpha_1\oplus\alpha_2\oplus\alpha_3\oplus V,
\tag{2.18}\label{eq:genus-two-decomposition}
\]
where \(V\) is the rank-four Schr\"odinger--Raynaud bundle.  The strict
rank theorem below forces
\(\min_p h^0(V(p))=1\), and this minimum holds on a nonempty open set.
For each nontrivial \(\alpha_i\), the curves \(C\) and
\(C-\alpha_i\) in \(J(C)\) have no common component, so
\(h^0(\alpha_i(p))=0\) for general \(p\).  Hence a general fibre has
exactly two sections.
This lowers the degree in the every-curve conclusion from \(36\) to
\(8\), compared with \cite{LandesmanLittPrill}.  This construction is
uniform over every genus-two curve, but it is not globally minimal.
The degree-five example in \cref{thm:C5-degree-five} shows that the
minimum over all genus-two curves and connected Prill-exceptional
covers is five.

The Galois closure gives a complementary sharp count.  Since the
regular representation of \(H_{2,2}\) decomposes as
\[
 \widetilde f_*\cO_{\widetilde X}
 \simeq
 \bigoplus_{\eta\in\widehat{(\bZ/2)^4}}\alpha_\eta
 \oplus V^{\oplus4},
\]
the degree-\(32\) cover has source genus \(33\), and every fibre has at least five
sections.  On the same minimum-open set used above, all nontrivial
\(\alpha_\eta(p)\) have no section, so a general fibre has exactly five.
\end{example}


\subsection{The surface scalar}

Write
\[
 \pi_g=\left\langle a_1,b_1,\ldots,a_g,b_g\ \middle|\
 \prod_{i=1}^g[a_i,b_i]=1\right\rangle.
\]

\begin{lemma}[Surface lifting obstruction]\label{lem:surface-lift}
Let \(A\) be a finite abelian group equipped with a projective
representation whose commutator is an alternating pairing
\(e\colon A\times A\to\bC^*\).  For a homomorphism
\(\alpha\colon\pi_g\to A\), a choice of lifts of
\(\alpha(a_i),\alpha(b_i)\) to the associated central extension satisfies
the surface relation if and only if
\[
 \kappa(\alpha):=
 \prod_{i=1}^g e\bigl(\alpha(a_i),\alpha(b_i)\bigr)=1.
\tag{2.19}\label{eq:surface-obstruction}
\]
The condition is independent of all chosen lifts.
\end{lemma}

\begin{proof}
Changing a lift by a central scalar does not change a commutator.  The
image of the defining surface relation is therefore exactly the product
in \eqref{eq:surface-obstruction}.  If the product is one, arbitrary
generator lifts define a homomorphism; if it is not one, no scalar
change can repair the relation.
\end{proof}

For the mod-\(n\) abelianization in a marked symplectic basis,
\(\kappa=\zeta^{-g}\).  Thus the standard projective Schr\"odinger
representation has a finite linear surface lift precisely when
\(n\mid g\).  Under this assumption define
\[
 \varphi_{g,n}\colon\pi_g\twoheadrightarrow H_{g,n},
 \qquad a_i\longmapsto x_i,\quad b_i\longmapsto y_i,
\tag{2.20}\label{eq:surface-heis-map}
\]
and set
\[
\rho_{g,n}=\tau_{g,n}\circ\varphi_{g,n}.
\tag{2.21}\label{eq:rho-definition}
\]
The map is surjective because the commutators generate the centre.

\section{Raynaud bundles and projective transport}

\subsection{The theta-equivariant Fourier--Mukai bundle}

Let \(C\) be a marked smooth projective curve of genus \(g\), let
\(J=\Pic^0(C)\), and choose an Abel--Jacobi embedding
\(\iota\colon C\hookrightarrow J\).  Choose a symmetric principal theta
divisor \(\Theta\), and put \(N=\cO_J(n\Theta)\).  If
\(\Phi_{\mathcal P}\) denotes the Fourier--Mukai transform for the
normalized Poincar\'e bundle, set
\[
 F_n=\Phi_{\mathcal P}(N^{-1})[g],\qquad R_n=\iota^*F_n.
\tag{3.1}\label{eq:raynaud-definition}
\]
Raynaud and Schneider prove the following facts
\cite{Raynaud,Schneider}:
\[
 \rk R_n=n^g,\qquad \mu(R_n)=\frac gn,\qquad
 H^0(C,R_n\otimes\alpha)\neq0
 \quad(\alpha\in\Pic^0(C)),
\tag{3.2}\label{eq:raynaud-support}
\]
and \(R_n\) is stable.  The key equivariant identity is
\[
 [n]^*F_n\simeq
 \cO_J(n\Theta)\otimes
 H^0(J,\cO_J(n\Theta))^\vee
\tag{3.3}\label{eq:schneider-splitting}
\]
in the following precise equivariant sense; see
\cite[Proposition 2.1(3)]{Schneider}.  Under the principal-polarization
identification \(J\simeq\widehat J\), the polarization isogeny of
\(\cO_J(n\Theta)\) is \([n]\).  Its theta group
\(\cG_n=\cG(\cO_J(n\Theta))\) acts on the right-hand side of
\eqref{eq:schneider-splitting}: the centre \(\bC^*\) has weight \(+1\)
on the line-bundle factor and weight \(-1\) on the dual section factor.
The tensor product therefore has central weight zero and carries an
honest action of \(\cG_n/\bC^*=J[n]\).  A theta structure identifies
the action on \(H^0(J,\cO_J(n\Theta))^\vee\) with the fixed geometric
Schr\"odinger module \(S\) in \eqref{eq:schrodinger}.  Its commutator is
the inverse Weil pairing, as prescribed by the chosen central-weight
convention.

Let
\[
 q\colon\widetilde C=C\times_{J,[n]}J\longrightarrow C.
\tag{3.4}\label{eq:kummer-cover}
\]
The cover is connected: its monodromy is the surjective mod-\(n\)
abelianization
\(\pi_1(C)\twoheadrightarrow H_1(C,\bZ/n)\).  Restricting
\eqref{eq:schneider-splitting} gives
\[
 q^*R_n\simeq\mathcal N\otimes S
\tag{3.5}\label{eq:raynaud-pullback}
\]
for a line bundle \(\mathcal N\) on \(\widetilde C\), with the full theta
group acting on the descent data.

\begin{lemma}[Projective theta descent]\label{lem:projective-theta}
After the marking identifies
\(J[n]\) with \(A_{g,n}\), the projective bundle \(\bP(R_n)\) is the
bundle associated to the fixed projective representation
\[
 \pi_g\longrightarrow A_{g,n}
 \longrightarrow PGL(S)
\tag{3.6}\label{eq:fixed-projective}
\]
given by the mod-\(n\) abelianization and the projective
Schr\"odinger representation on \(S\).  Its projective conjugacy class
is independent of \(C\), of
the Abel--Jacobi base point, and of the chosen theta structure.
\end{lemma}

\begin{proof}
Projectivizing \eqref{eq:raynaud-pullback} removes the line factor
\(\mathcal N\), so \(q^*\bP(R_n)\) is the constant
\(\bP(S)\)-bundle.  Write \(G=J[n]\) for the deck group.  Under the
theta-equivariant form of \eqref{eq:schneider-splitting}, the descent
isomorphism attached to \(a\in G\) is represented on
\(\mathcal N\otimes S\) by
\[
 \ell_a\otimes T_a,
\]
where \(\ell_a\colon a^*\mathcal N\to\mathcal N\) is a line isomorphism
and \(T_a\in GL(S)\) is a theta-group lift.  Replacing a lift by a scalar
rescales \(\ell_a\) inversely, so \([T_a]\in PGL(S)\) is intrinsic.  The
descent cocycle gives
\([T_{a+b}]=[T_aT_b]\), while
\[
 T_aT_bT_a^{-1}T_b^{-1}=e_n(a,b)I.
\]
Thus the projective deck action is the constant Schr\"odinger action
with commutator the inverse Weil pairing.  The cover \(q\) is the
pullback of \([n]\) along the Abel--Jacobi map, hence its monodromy is
the marked mod-\(n\) abelianization.  The marking and the convention in
\eqref{eq:schrodinger} therefore identify the descent representation
with \eqref{eq:fixed-projective}.

Choose invariant Hermitian metrics on the finite theta representation
and on the descended line factor.  The preceding cocycle then defines a
unitary flat connection on \(\bP(R_n)\) whose holomorphic structure is
the one obtained from theta descent.  The stable bundle \(R_n\) has a
unique Hermitian--Einstein connection up to unitary scalar gauge, and its
projectivization is flat.  Uniqueness identifies its projective
connection with the descended unitary connection.  Thus algebraic theta
descent and projective Narasimhan--Seshadri monodromy give the same
projective local system.

Changing the theta structure conjugates the projective representation.
Changing the theta divisor or the Abel--Jacobi base point translates a
semihomogeneous Fourier--Mukai bundle; the translation formula
\(t_a^*F_n\simeq F_n\otimes P_b\), for a suitable
\(P_b\in\Pic^0(J)\), changes its restriction only by a line twist.
Consequently the projectivization is unchanged.  There is no even-level
defect: since \(\bC^*\) is divisible,
\[
 H^2(A_{g,n},\bC^*)\simeq
 \Hom(\wedge^2A_{g,n},\bC^*),
\]
so the projective cohomology class is determined by its commutator.
Quadratic refinements affect extensions of symplectic automorphisms, not
the marked projective representation used here.
\end{proof}

\subsection{Scalar transport of support loci}

For a bundle \(E\) on \(C\), define
\[
 V_d^i(E)=
 \{L\in\Pic^d(C):H^i(C,E\otimes L)\neq0\}.
\tag{3.7}\label{eq:support-locus}
\]

\begin{lemma}[Projective scalar transport]\label{lem:scalar-transport}
Let \(E\) and \(E'\) be vector bundles of the same rank on a curve.  If
\(\bP(E)\simeq\bP(E')\) over the curve, then
\[
 E'\simeq E\otimes\alpha
 \quad\text{for some line bundle }\alpha.
\tag{3.8}\label{eq:scalar-twist}
\]
If \(\deg E=\deg E'\), then \(\alpha\in\Pic^0(C)\).
Consequently,
\[
 V_d^i(E')=\alpha^{-1}V_d^i(E).
\tag{3.9}\label{eq:support-translate}
\]
In particular, the property \(V_d^i(E)=\Pic^d(C)\) depends only on the
projective local system.
\end{lemma}

\begin{proof}
The two bundles define the same \(PGL_r\)-torsor.  Comparing lifts through
\(1\to\bG_m\to GL_r\to PGL_r\to1\) shows that their transition matrices
differ by a \(\bG_m\)-cocycle, which is a line bundle \(\alpha\).  This
proves \eqref{eq:scalar-twist}.  Taking degrees gives
\(\deg E'=\deg E+r\deg\alpha\), so equal degrees force
\(\deg\alpha=0\).  Formula \eqref{eq:support-translate} follows by
tensoring line bundles.
\end{proof}

\begin{theorem}[Heisenberg--Raynaud uniformization]\label{thm:uniform}
Let \(g\geq2\), \(n\geq2\), and \(n\mid g\).  The representation
\(\rho_{g,n}\) in \eqref{eq:rho-definition} satisfies all conclusions
of \cref{thm:intro-uniform}.
\end{theorem}

\begin{proof}
Put \(m=g/n\) and \(r=n^g\).  Choose
\(M\in\Pic^{-m}(C)\) with
\[
 M^{\otimes r}\simeq(\det R_n)^{-1};
\tag{3.10}\label{eq:det-normalization}
\]
such an \(M\) exists because multiplication by \(r\) is surjective on
\(\Pic^0(C)\).  The bundle
\[
 W=R_n\otimes M
\]
is stable of degree zero and has trivial determinant.  By the
Narasimhan--Seshadri theorem \cite{NarasimhanSeshadri}, it is associated
to an irreducible unitary representation \(\sigma_C\).

The projective local system of \(\sigma_C\) is the one in
\cref{lem:projective-theta}.  The fixed finite representation
\(\rho_{g,n}\) is a linear lift of the same projective representation by
\cref{lem:surface-lift}.  Therefore
\cref{lem:scalar-transport} gives an \(\alpha_C\in\Pic^0(C)\) such that
\[
 \cV_{\rho_{g,n},C}
 \simeq R_n\otimes M\otimes\alpha_C.
\tag{3.11}\label{eq:uniform-identification}
\]
For \(L\in\Pic^m(C)\), the line
\(M\otimes\alpha_C\otimes L\) has degree zero.  Equation
\eqref{eq:raynaud-support} now proves
\[
 H^0(C,\cV_{\rho_{g,n},C}\otimes L)\neq0.
\]

The representation is irreducible and nontrivial, so its degree-zero
flat bundle is stable and has no global sections.  Triviality of the
determinant follows directly from \cref{lem:character}.  Its image is
the faithful Schr\"odinger image of \(H_{g,n}\), hence has order
\(n^{2g+1}\).
\end{proof}

\begin{remark}[Logical scope]\label{rem:scope-uniform}
The uniformity in \cref{thm:uniform} is not a compactness or
finite-candidate argument.  It uses two exact inputs:
theta-equivariant projective descent and invariance of the full Picard
support condition under scalar twists.  A statement only along the Abel
curve \(p\mapsto\cO_C(p)\) would not be sufficient for this transport.
\end{remark}

\section{Lagrangian descent and Heisenberg untwisting}

\subsection{The three-step factorization}

Put
\[
 K=\langle y_1,\ldots,y_g\rangle,\qquad
 B=\langle z\rangle K.
\tag{4.1}\label{eq:lagrangian-subgroups}
\]
Then \(K\simeq(\bZ/n)^g\), \(B=\langle z\rangle\times K\), and
\[
 |K|=r,\qquad |B|=nr.
\tag{4.2}\label{eq:BK-orders}
\]
The subgroup \(B\) is normal in \(H_{g,n}\), since conjugation changes
an element of \(B\) only by a central element.  Define
\[
 \chi\colon B\longrightarrow\mu_n,\qquad
 \chi(z)=\zeta^{-1},\quad \chi(y_i)=1.
\tag{4.3}\label{eq:chi}
\]
The line \(\bC e_0\subset S\) is a \(B\)-eigenline of character \(\chi\).
It is cyclic under the shifts, and therefore
\[
 S\simeq\Ind_B^{H_{g,n}}\chi.
\tag{4.4}\label{eq:monomial-schrodinger}
\]

\begin{theorem}[Exact untwisting]\label{thm:exact-untwisting}
Let \(C\) be marked and retain the hypotheses of \cref{thm:uniform}.
Let
\[
 \pi_B=\varphi_{g,n}^{-1}(B),\qquad
 \pi_K=\varphi_{g,n}^{-1}(K).
\]
The corresponding covers form a diagram
\[
 X\xrightarrow{\,v\,}Y\xrightarrow{\,u\,}C
\tag{4.5}\label{eq:untwisting-diagram}
\]
with the following properties.
\begin{enumerate}[label=\textup{(\alph*)}]
\item The cover \(u\) is connected Galois of degree \(r\).
\item The character \(\chi\circ\varphi_{g,n}|_{\pi_B}\) defines an
exact-order-\(n\) torsion line bundle \(\lambda\) on \(Y\), chosen with
the associated-bundle convention for which
\[
 u_*\lambda\simeq\cV_{\rho_{g,n},C}.
\tag{4.6}\label{eq:push-line}
\]
If the opposite convention is used, \(\lambda\) is replaced by
\(\lambda^{-1}\).
\item The cover \(v\) is connected cyclic \'etale of degree \(n\),
\(v^*\lambda\simeq\cO_X\), and
\[
 X\simeq
 \Spec_Y\!\left(\bigoplus_{j=0}^{n-1}\lambda^{-j}\right)
\tag{4.7}\label{eq:cyclic-spec}
\]
after a choice of \(\lambda^n\simeq\cO_Y\).
\item For every line bundle \(L\) on \(C\) and \(i=0,1\),
\[
 H^i(Y,\lambda\otimes u^*L)
 \simeq H^i(C,\cV_{\rho_{g,n},C}\otimes L).
\tag{4.8}\label{eq:all-cohomology-line}
\]
\item The composite \(f=u\circ v\) has degree \(nr=n^{g+1}\), and
\[
 f_*\cO_X\simeq
 \cO_C\oplus\cV_{\rho_{g,n},C}\oplus E
\tag{4.9}\label{eq:ordinary-splitting}
\]
for a unitary degree-zero bundle \(E\).
\item The Galois closure of \(f\) is the \(H_{g,n}\)-cover and has degree
\(nr^2=n^{2g+1}\).
\end{enumerate}
\end{theorem}

\begin{proof}
Equation \eqref{eq:monomial-schrodinger} and the equivalence between
finite-cover direct image and induction give \eqref{eq:push-line}.  The
character is surjective because its restriction to the centre is
faithful; hence \(\lambda\) has exact order \(n\).  Its kernel in
\(\pi_B\) is exactly \(\pi_K\), which proves the assertions about \(v\)
and \eqref{eq:cyclic-spec}.  Finite direct image is exact, and the
projection formula proves \eqref{eq:all-cohomology-line}.

The permutation module of the composite is
\(\bC[H_{g,n}/K]=\Ind_K^{H_{g,n}}\mathbf1\).  Frobenius reciprocity and
the fixed lines in both \(S\) and \(S^\vee\) show that it contains
\(S\) once.  It also contains the trivial representation, proving
\eqref{eq:ordinary-splitting}.

It remains to determine the Galois closure.  If \(k\) belongs to the core
of \(K\), then \([h,k]\in K\cap\langle z\rangle=1\) for every \(h\in H\).
The nondegeneracy of \eqref{eq:weil-pairing} forces the image of \(k\) in
\(H/\langle z\rangle\) to vanish, and then \(k=1\).  Thus the core is
trivial.
\end{proof}

\subsection{Finite-symplectic theta untwisting}

The preceding argument is not restricted to a homogeneous module
\((\bZ/n)^{2g}\).  The following form is the group-theoretic engine
needed for arbitrary polarization type.

\begin{theorem}[Finite-symplectic theta untwisting]
\label{thm:finite-symplectic-untwisting}
Let \((A,e)\) be a finite abelian group with a perfect alternating
pairing \(e\colon A\times A\to\mu_N\), where \(N\) is the exponent of
\(A\), assume \(N\geq2\), and write \(|A|=r^2\).  Choose a Lagrangian
decomposition \(A=L\oplus L'\), and let \(H_{\mathrm{sp}}(A)\) be the
split finite Heisenberg model described in the proof.  Suppose a surface group admits a
surjection
\(\psi\colon\pi_1(C)\twoheadrightarrow H_{\mathrm{sp}}(A)\).  For a faithful
central character, let \(S_A\) be the corresponding Schr\"odinger
module and \(V_A\) its flat bundle.  Then there are connected covers
\[
 X_A\longrightarrow Y_A\longrightarrow C
\]
and an exact-order-\(N\) line bundle \(\lambda_A\in\Pic^0(Y_A)\) such
that
\[
 \deg(Y_A/C)=r,\qquad
 (u_A)_*\lambda_A\simeq V_A,
\]
the first map is the cyclic degree-\(N\) torsor trivializing
\(\lambda_A\), and
\[
 (f_A)_*\cO_{X_A}\simeq\cO_C\oplus V_A\oplus E_A.
\]
Thus \(\deg f_A=Nr\), and the Galois closure of \(f_A\) is the
\(H_{\mathrm{sp}}(A)\)-cover of degree \(Nr^2\).  For \(i=0,1\), finite direct image
identifies the cohomology of \(V_A\otimes L\) with that of
\(\lambda_A\otimes u_A^*L\).
\end{theorem}

\begin{proof}
The elementary-divisor theorem for finite symplectic modules gives a
decomposition
\[
 A=L\oplus L',\qquad |L|=|L'|=r,
\]
in which both summands are isotropic and the pairing between them is
perfect.  Concretely, if
\(A\simeq\bigoplus_i(\bZ/d_i)^2\), take one cyclic factor from each
  hyperbolic plane.  The split finite model in the statement is
\[
 H_{\mathrm{sp}}(A)=\mu_N\times L\times L',
\]
with multiplication
\[
 (z,l,l')(w,m,m')=
 \bigl(zw\,e(l',m),\ l+m,\ l'+m'\bigr).
\]
Its commutator is the prescribed pairing.  The projective cocycle class
with commutator \(e\) is unique because
\[
 H^2(A,\bC^*)\simeq\Hom(\wedge^2A,\bC^*),
\]
and the finite Stone--von Neumann theorem gives the unique irreducible
Schr\"odinger projective representation for each faithful central
character.  This does not assert that every finite central extension
with the same commutator is isomorphic to the split group.

Put \(B=\mu_N\times L\).  For a faithful character
\(\omega\colon\mu_N\to\bC^*\), extend it to
\(\chi(z,l)=\omega(z)\) on \(B\).  Then
\[
 S_A\simeq\Ind_B^{H_{\mathrm{sp}}(A)}\chi,
\]
so \(\dim S_A=[H_{\mathrm{sp}}(A):B]=r\).  The group \(B\) is normal.  The subgroups
\(\psi^{-1}(B)\) and \(\psi^{-1}(L)\) therefore define the asserted
degree-\(r\) cover and its cyclic degree-\(N\) refinement.  Surjectivity
of \(\psi\) and faithfulness of \(\omega\) show that the associated
character line has exact order \(N\).  Induction, finite direct image,
and the projection formula give the direct-image and cohomology
identities.

The permutation module \(\bC[H_{\mathrm{sp}}(A)/L]\) contains the trivial module and
also \(S_A\), since the dual Schr\"odinger module has an \(L\)-fixed
inducing line and Frobenius reciprocity applies.  Finally, the core of
\(L\) is trivial.  Indeed, if
\(l\in L\) belongs to every conjugate of \(L\), conjugation by every
\(l'\in L'\) forces \(e(l',l)=1\); perfectness gives \(l=0\).  Hence the
Galois closure has group \(H_{\mathrm{sp}}(A)\) and order \(Nr^2\).
\end{proof}

\begin{corollary}[Higher-divisor Prill condition]\label{cor:higher-prill}
Let \(f=u\circ v\colon X\to C\) be the composite in
\cref{thm:exact-untwisting}.  Then, for every effective divisor \(D\) of
degree \(m=g/n\),
\[
 h^0\!\left(X,\cO_X(f^*D)\right)\geq2.
\tag{4.10}\label{eq:higher-prill}
\]
More strongly, for every \(L\in\Pic^m(C)\),
\[
 h^0(X,f^*L)\geq h^0(C,L)+1.
\tag{4.11}\label{eq:full-pic-pullback}
\]
\end{corollary}

\begin{proof}
Apply \eqref{eq:ordinary-splitting}, the projection formula, and
\cref{thm:uniform}.  If \(D\) is effective, the trivial summand gives one
section and the Schr\"odinger summand gives another.
\end{proof}

\section{Brill--Noether geometry of the compression tower}

Retain \(m=g/n\), and let \(f_k\colon X_k\to C\) be the covers from
\cref{prop:compression-tower}.  Put
\[
 G_k=g(X_k)=1+d_k(g-1),\qquad e_k=md_k.
\tag{5.1}\label{eq:BN-numbers}
\]
I use the convention
\[
 W_e^s(X)=
 \{L\in\Pic^e(X):h^0(X,L)\geq s+1\}.
\]

\begin{theorem}[Abelian and symmetric-product blocks]\label{thm:BN-blocks}
For every \(0\leq k\leq g\),
\[
 f_k^*\Pic^m(C)\subset W_{e_k}^{\mu_k-1}(X_k)
\tag{5.2}\label{eq:full-pic-BN}
\]
is a translate of an abelian subvariety of dimension \(g\).  Moreover,
\[
 f_k^*W_m^0(C)\subset W_{e_k}^{\mu_k}(X_k)
\tag{5.3}\label{eq:effective-BN}
\]
contains the image of an \(m\)-dimensional symmetric product.  The
corresponding Brill--Noether numbers are
\[
 \rho_k=
 1+d_k(g-1)
 -\mu_k\bigl[d_k(g-1-m)+\mu_k\bigr]
\tag{5.4}\label{eq:rho-full}
\]
and
\[
 \rho_k^+=
 1+d_k(g-1)
 -(\mu_k+1)
 \bigl[1+d_k(g-1-m)+\mu_k\bigr].
\tag{5.5}\label{eq:rho-effective}
\]
\end{theorem}

\begin{proof}
Equation \eqref{eq:tower-sections} gives
\(h^0(X_k,f_k^*L)\geq\mu_k\) for every \(L\in\Pic^m(C)\), and gives
\(h^0(X_k,f_k^*L)\geq\mu_k+1\) when \(L\) is effective.  This proves the
two inclusions, where \(C^{(m)}\to W_m^0(C)\) is the Abel--Jacobi map.

The identity
\[
 \Nm_{f_k}\circ f_k^*=[d_k]\colon\Pic^0(C)\longrightarrow\Pic^0(C)
\tag{5.6}\label{eq:norm-pull}
\]
shows that \(f_k^*\) has finite kernel.  Hence the first image has
dimension \(g\).  Since \(m\leq g-1\), a general effective divisor of
degree \(m\) has one section, so \(\dim W_m^0(C)=m\); the constructed
second image therefore has dimension \(m\).  Substitution of
\(G_k=1+d_k(g-1)\), \(e_k=md_k\), and the relevant values of \(s\) in
\[
 \rho(G,s,e)=G-(s+1)(G-e+s)
\]
gives \eqref{eq:rho-full} and \eqref{eq:rho-effective}.
\end{proof}

\begin{corollary}[Negative expected dimension]\label{cor:negative-BN}
At the minimal ordinary layer \(k=g\), put \(d=n^{g+1}\) and write
\(f=f_g\colon X=X_g\to C\).  Then
\[
 f^*\Pic^m(C)\subset W_{dm}^0(X),\qquad
 f^*C^{(m)}\subset W_{dm}^1(X).
\tag{5.7}\label{eq:minimal-BN}
\]
The second locus has
\[
 \rho\bigl(g(X),1,dm\bigr)=d(2m-g+1)-3.
\tag{5.8}\label{eq:minimal-rho}
\]
If \(n\geq3\), this number is negative, although the constructed image
\(f^*C^{(m)}\) has dimension \(m\).  In particular, for \(n=g\)
and \(g\geq3\),
\[
 f^*C\subset W_d^1(X),\qquad
 \rho\bigl(g(X),1,d\bigr)=d(3-g)-3<0.
\tag{5.9}\label{eq:Prill-negative-rho}
\]
\end{corollary}

\begin{proof}
Here \(\mu_g=1\).  Formula \eqref{eq:minimal-rho} follows directly from
\(g(X)=1+d(g-1)\).  If \(n\geq3\), then \(g\geq n\) and
\[
 2m-g+1=\frac{2g}{n}-g+1\leq0;
\]
equality occurs only at the boundary \(g=n=3\), where the remaining
term is \(-3\).
\end{proof}

Write
\[
 d_s(X)=\min\{\deg L:L\in\Pic(X),\ h^0(X,L)\geq s+1\}
\]
for the gonality sequence.

\begin{corollary}[Gonality sequence and Clifford index]\label{cor:clifford-tower}
For every effective divisor \(D\) of degree \(m\),
\[
 h^0(X_k,\cO_{X_k}(f_k^*D))\geq\mu_k+1,\qquad
 h^1(X_k,\cO_{X_k}(f_k^*D))
 \geq\mu_k+1+d_k(g-1-m).
\tag{5.10}\label{eq:special-pullback}
\]
Consequently,
\[
 d_{\mu_k}(X_k)\leq e_k,\qquad
 \Cliff(X_k)\leq e_k-2\mu_k.
\tag{5.11}\label{eq:gonality-clifford}
\]
\end{corollary}

\begin{proof}
The first inequality is \eqref{eq:tower-sections}.  Riemann--Roch gives
the second.  Since \(m\leq g/2\), both \(h^0\) and \(h^1\) are at least
two, and the definitions of the gonality sequence and Clifford index
give \eqref{eq:gonality-clifford}.
\end{proof}

\section{A strict rank--corank inequality}

The argument in this section applies to arbitrary stable degree-zero
bundles; neither finite monodromy nor the theta construction is used.

\begin{theorem}[Strict rank--corank inequality]\label{thm:strict-rank}
Let \(C\) have genus \(g\geq2\), and let \(W\) be a nontrivial stable
bundle of degree zero and rank \(r\).  Set
\[
 c=\min_{p\in C}h^0(C,W(p)).
\]
If \(c\geq1\), then
\[
 r\geq gc+2.
\tag{6.1}\label{eq:strict-rank}
\]
\end{theorem}

\begin{proof}
Set \(F=K_C\otimes W^\vee\).  Then \(F\) is stable of rank \(r\), degree
\(2r(g-1)\), and
\[
 h^0(F)=r(g-1),\qquad h^1(F)=h^0(W)=0.
\tag{6.2}\label{eq:F-sections}
\]
Let \(U\subset F\) be the saturation of the image of the evaluation map
\[
 H^0(F)\otimes\cO_C\longrightarrow F,
\]
and put \(s=\rk U\), \(D=\deg U\).  At a point where \(h^0(W(p))=c\),
the cokernel of evaluation in the fibre of \(F\) has dimension \(c\):
this follows from the exact sequence for \(F(-p)\), because
\[
 H^1(F(-p))^\vee\simeq H^0(W(p)).
\]
Consequently,
\[
 s=r-c,\qquad h^0(U)=h^0(F)=r(g-1).
\tag{6.3}\label{eq:U-data}
\]
Indeed, every global section of \(F\) lands generically in the
evaluation image and therefore factors through its saturation \(U\);
the reverse inclusion of section spaces is tautological.

The bundle \(U\) is generically globally generated.  Its minimal
Harder--Narasimhan quotient is therefore generically globally generated
and has nonnegative degree; hence every Harder--Narasimhan slope of
\(U\) is nonnegative.  Its maximal-slope subbundle is a proper subbundle
of the stable bundle \(F\), so every such slope is strictly less than
\(2g-2\).  The vector-bundle Clifford inequality
\cite[Theorem 2.1]{BrambilaPazGN}, applied to the semistable
Harder--Narasimhan factors and then added, gives
\[
 h^0(U)\leq s+\frac D2.
\]
For a slope-zero factor, the same bound follows directly: generic global
generation and degree zero make it trivial, so its number of sections is
its rank.  Thus no endpoint outside the cited positive-slope range is
being used.
Together with \eqref{eq:U-data}, this yields
\[
 D\geq2r(g-2)+2c.
\tag{6.4}\label{eq:D-lower}
\]
Stability of \(F\), integrality of \(D\), and \(s<r\) give
\[
 D\leq2s(g-1)-1.
\tag{6.5}\label{eq:D-upper}
\]
Equations \eqref{eq:D-lower} and \eqref{eq:D-upper} first imply
\(r\geq gc+1\).

It remains to exclude the endpoint \(r=gc+1\).  Then
\[
 s=(g-1)c+1,
\]
and the two bounds for \(D\) differ by one.  Thus
\[
 D=2s(g-1)-e,\qquad e\in\{1,2\}.
\tag{6.6}\label{eq:D-endpoint}
\]
Define
\[
 E=K_C\otimes U^\vee.
\]
Riemann--Roch, Serre duality, \eqref{eq:U-data}, and
\eqref{eq:D-endpoint} give
\[
 \rk E=s,\qquad \deg E=e,\qquad h^0(E)=s+e-1.
\tag{6.7}\label{eq:E-data}
\]
Writing \(Q=F/U\), saturation makes \(Q\) locally free.  Dualizing and
tensoring by \(K_C\) gives a surjection
\[
 W\twoheadrightarrow E.
\tag{6.8}\label{eq:W-to-E}
\]
Since \(\rk E=s=r-c<r\), every nonzero vector-bundle quotient of \(E\),
composed with \eqref{eq:W-to-E}, is a proper vector-bundle quotient of
the stable degree-zero bundle \(W\), and consequently has strictly
positive degree.

I next show that \(E\) is generically globally generated.  Let
\(A\subset E\) be the saturation of its evaluation image, and suppose
\(a=\rk A<s\).  The quotient \(E/A\) has positive degree, so
\[
 0\leq\deg A\leq e-1.
\]
All sections of \(E\) factor through \(A\), hence \(h^0(A)=h^0(E)\).
Choosing \(a\) generically independent sections of the generically
generated bundle \(A\) gives
\[
 0\longrightarrow\cO_C^a\longrightarrow A\longrightarrow T_A
 \longrightarrow0,
\]
where \(T_A\) has length \(\deg A\).  It follows that
\[
 h^0(E)=h^0(A)\leq a+\deg A
 \leq s+e-2,
\]
contrary to \eqref{eq:E-data}.  Thus \(E\) is generically globally
generated.

Choose \(s\) generically independent global sections.  Their determinant
gives
\[
 0\longrightarrow\cO_C^s\longrightarrow E\longrightarrow T
 \longrightarrow0,\qquad \operatorname{length}T=e.
\tag{6.9}\label{eq:E-evaluation}
\]
Dualizing gives
\[
 0\longrightarrow E^\vee\longrightarrow\cO_C^s
 \longrightarrow\mathcal Ext^1(T,\cO_C)\longrightarrow0.
\tag{6.10}\label{eq:E-dual-evaluation}
\]
If \(s>e\), the map
\(\bC^s\to H^0(\mathcal Ext^1(T,\cO_C))\), whose target has dimension
\(e\), has a nonzero kernel.  The resulting section of \(E^\vee\) maps
to a nonzero constant vector in \(\cO_C^s\), so it vanishes nowhere.
Dualizing this line subbundle gives a quotient \(E\twoheadrightarrow
\cO_C\), contradicting the positivity of all nonzero quotients of \(E\).
Therefore \(s\leq e\).

Since \(s=(g-1)c+1\geq2\) and \(e\leq2\), the only remaining case is
\[
 s=e=2,\qquad g=2,\qquad c=1.
\tag{6.11}\label{eq:last-endpoint}
\]
Now \(E\) has rank two, degree two, and three global sections.  The
wedge map
\[
 \wedge^2H^0(E)\longrightarrow H^0(\det E)
\tag{6.12}\label{eq:wedge-map}
\]
is injective.  Indeed, every nonzero bivector in
\(\wedge^2\bC^3\) is decomposable; a nonzero element of the kernel would
give two independent sections contained in a line subbundle
\(L\subset E\).  The quotient \(E/L\) has positive degree, so
\(\deg L\leq1\), whereas \(h^0(L)\geq2\), impossible on a genus-two
curve.  Injectivity in \eqref{eq:wedge-map} would imply
\(h^0(\det E)\geq3\).  A degree-two line bundle on a genus-two curve has
at most two sections.  This final contradiction excludes
\eqref{eq:last-endpoint} and proves \eqref{eq:strict-rank}.
\end{proof}

\begin{corollary}[Degree and section bounds for covers]\label{cor:cover-lower-bound}
Let \(f\colon X\to C\) be a connected finite \'etale cover of a
genus-\(g\) curve.  If
\[
 \min_{p\in C}
 h^0\!\left(X,\cO_X(f^{-1}(p))\right)\geq k+1
\quad(k\geq1),
\tag{6.13}\label{eq:k-sections}
\]
then
\[
 \deg f\geq gk+3.
\tag{6.14}\label{eq:degree-lower}
\]
In particular, every connected Prill-exceptional \'etale cover has
degree at least \(g+3\).
\end{corollary}

\begin{proof}
Decompose the unitary permutation bundle as
\[
 f_*\cO_X\simeq
 \cO_C\oplus\bigoplus_j W_j^{\oplus m_j},
\tag{6.15}\label{eq:permutation-decomp}
\]
where the \(W_j\) are nontrivial stable degree-zero bundles.  Put
\[
 c_j=\min_p h^0(W_j(p)).
\]
The minimum of each upper-semicontinuous function is attained on a
nonempty open subset.  On the intersection of these finitely many open
sets, \eqref{eq:k-sections} gives
\[
 \sum_jm_jc_j\geq k.
\]
For every \(j\) with \(c_j>0\), \cref{thm:strict-rank} gives
\(\rk W_j\geq gc_j+2\).  Hence
\[
 \deg f
 =1+\sum_jm_j\rk W_j
 \geq1+gk+2\sum_{c_j>0}m_j
 \geq gk+3.
\]
\end{proof}

\begin{corollary}[Galois lower bound]\label{cor:galois-lower}
If a connected Galois \'etale cover \(f\colon X\to C\) is
Prill-exceptional, then its deck group \(G\) satisfies
\[
 |G|\geq(g+2)(g+3).
\tag{6.16}\label{eq:galois-order-lower}
\]
\end{corollary}

\begin{proof}
Write the regular decomposition as
\[
 f_*\cO_X\simeq\cO_C\oplus
 \bigoplus_{\rho\neq\mathbf1}W_\rho^{\oplus r_\rho},
 \qquad r_\rho=\rk W_\rho.
\]
On the common nonempty open set where every
\(h^0(W_\rho(p))\) attains its minimum \(c_\rho\),
Prill-exceptionality gives
\(\sum_{\rho\neq\mathbf1}r_\rho c_\rho\geq1\).  Hence some nontrivial
summand \(W=W_\rho\) has \(c=c_\rho>0\).  If \(r=\rk W\), then
\(r\geq g+2\) by \cref{thm:strict-rank}.  The regular representation
contains the trivial representation and \(W\) with multiplicity \(r\),
so \(|G|>r^2\).  The degree of a complex irreducible character divides
the group order by the Frobenius divisibility theorem
\cite[\S6.5]{SerreFiniteGroups}; hence \(|G|\) is a
multiple of \(r\), and
\[
 |G|\geq r(r+1)\geq(g+2)(g+3).
\]
\end{proof}

\begin{corollary}[A numerical Putman--Wieland range]\label{cor:PW-range}
A connected finite unramified cover of a closed genus-\(g\) surface with
degree at most \(g+2\) cannot furnish a counterexample to the
Putman--Wieland conjecture \cite{PutmanWieland}.  If the cover is Galois, the same conclusion
holds whenever the deck group has order less than \((g+2)(g+3)\).
\end{corollary}

\begin{proof}
Landesman and Litt prove that a closed Putman--Wieland counterexample is
Prill-exceptional \cite[Lemma 5.5]{LandesmanLittApplications}.  Apply
\cref{cor:cover-lower-bound,cor:galois-lower}.
\end{proof}

\begin{remark}[No converse]\label{rem:no-PW-converse}
Prill-exceptionality is only a necessary condition for a
Putman--Wieland counterexample.  The Heisenberg covers constructed here
therefore do not by themselves produce finite homology orbits.  The
large-monodromy theorem of Landesman--Litt--Sawin
\cite{LandesmanLittSawin} applies in a complementary range in which the
base genus is large compared with the representation rank; the ranks
\(n^g\) considered here lie far outside that range.
\end{remark}

\subsection{The degree-five minimum in genus two}

The numerical lower bound leaves degree five as the first possible
value in genus two.  I first isolate its monodromy and theta
obstruction, and then construct a symmetric example for which the
obstruction vanishes.

\begin{lemma}[Rank-four no-jump lemma]\label{lem:rank-four-no-jump}
Let \(C\) have genus two, and let \(W\) be a nontrivial stable
degree-zero bundle of rank four.  If
\[
 \min_{p\in C}h^0(C,W(p))\geq1,
\]
then
\[
 h^0(C,W(p))=1\qquad\text{for every }p\in C.
\tag{6.17}\label{eq:rank-four-no-jump}
\]
\end{lemma}

\begin{proof}
The strict rank--corank inequality gives
\(4\geq2c+2\), so the minimum is \(c=1\).  Put
\[
 F=K_C\otimes W^\vee,\qquad V=H^0(C,F).
\]
Then \(\dim V=4\), and the evaluation map
\(V\otimes\cO_C\to F\) has generic rank three.  Its kernel \(M\) is a
line subbundle; set \(A=M^\vee\).  Let \(V_0\subset V\) be the smallest
constant subspace such that
\(M\subset V_0\otimes\cO_C\), and put \(t=\dim V_0\).  Minimality gives
an injection \(V_0^\vee\hookrightarrow H^0(C,A)\), and these sections
generate \(A\).  If \(I_0\) is the image of
\(V_0\otimes\cO_C\to F\) and \(U_0\) its saturation, then
\[
 \deg A=\deg I_0\leq\deg U_0\leq2(t-1)-1,
\tag{6.18}\label{eq:A-upper}
\]
because \(F\) is stable of slope two.  The case \(t=1\) would give a
constant linear relation among the elements of \(V=H^0(F)\), so
\(t\geq2\).

On a genus-two curve, a globally generated nontrivial line bundle with
at least \(t\) sections has degree at least \(2,4,5\) for
\(t=2,3,4\), respectively.  This contradicts
\eqref{eq:A-upper} for \(t=2,3\), so \(t=4\) and \(\deg A\geq5\).
Let \(I\) be the full evaluation image and \(U\) its saturation.  Then
\(\rk U=3\), stability gives \(\deg U\leq5\), and
\[
 \deg U=\deg I+\operatorname{length}(U/I)
       =\deg A+\operatorname{length}(U/I)\geq5.
\]
Consequently \(\deg A=\deg U=5\) and \(U=I\).  Evaluation is therefore
surjective onto \(U\) at every point, and its cokernel in \(F_p\) has
dimension one for every \(p\).  Serre duality identifies this cokernel
with \(H^0(C,W(p))^\vee\), proving
\eqref{eq:rank-four-no-jump}.
\end{proof}

\begin{proposition}[Degree-five monodromy reduction]
\label{prop:degree-five-reduction}
Let \(f\colon X\to C\) be a connected degree-five \'etale cover of a
genus-two curve.  If every fibre moves in a pencil, then the monodromy
group is \(A_5\) in its natural five-point action.  The augmentation
bundle \(W\) is a stable oriented orthogonal bundle of rank four and
satisfies
\[
 h^0(C,W(p))=1\quad(p\in C),\qquad
 w_1(W)=0,\qquad \langle w_2(W),[C]\rangle=1.
\tag{6.19}\label{eq:A5-topological-data}
\]
\end{proposition}

\begin{proof}
Decompose the augmentation local system into irreducible unitary
summands.  On the common open set where all section dimensions are
minimal, Prill-exceptionality forces one summand \(E\) to satisfy
\(\min_p h^0(E(p))>0\).  By \cref{thm:strict-rank},
\(\rk E\geq4\).  Since the full augmentation rank is four, the
augmentation representation is irreducible.  It is therefore stable,
and \cref{lem:rank-four-no-jump} gives the first assertion in
\eqref{eq:A5-topological-data}.

An irreducible augmentation representation is equivalent to
two-transitivity of the five-point action: the number of orbits on
ordered pairs is
\(\dim\operatorname{End}_G(\bC^5)\).  The two-transitive subgroups of
\(S_5\) are
\[
 \operatorname{AGL}_1(\mathbf F_5),\qquad A_5,\qquad S_5.
\tag{6.20}\label{eq:degree-five-groups}
\]
For completeness, two-transitivity makes \(|G|\) divisible by
\(5\cdot4\).  If \(|G|=20\), the action is sharply two-transitive and
\(G=C_5\rtimes C_4\).  Subgroups of orders \(40\) and \(60\) not equal
to \(A_5\) are excluded by intersecting with the simple group \(A_5\),
and order \(120\) gives \(S_5\).

It remains to use the orthogonal structure.  For a real orthogonal
bundle \(E_{\mathbf R}\) of rank \(r\) and a theta characteristic
\(\kappa\), the mod-two index formula \cite{AtiyahSpin} is
\[
 h^0(C,E_{\mathbf R}\otimes_{\mathbf R}\bC\otimes\kappa)
 \equiv r h^0(C,\kappa)+q_\kappa(w_1(E_{\mathbf R}))
 +\langle w_2(E_{\mathbf R}),[C]\rangle\pmod2,
\tag{6.21}\label{eq:nonoriented-parity}
\]
where \(q_\kappa\) is the Riemann--Mumford quadratic refinement.  One
may derive \eqref{eq:nonoriented-parity} from the mod-two index theorem
by the splitting principle: for real line classes \(a_i\), use
\(q_\kappa(\sum a_i)=\sum q_\kappa(a_i)+
\sum_{i<j}a_i\smile a_j\).

The six Weierstrass points give the six odd theta characteristics
\(\kappa_i=\cO_C(w_i)\).  Equation
\eqref{eq:rank-four-no-jump} and \eqref{eq:nonoriented-parity} imply
\[
 q_{\kappa_i}(w_1(W))+\langle w_2(W),[C]\rangle=1
 \qquad(1\leq i\leq6).
\tag{6.22}\label{eq:six-parities}
\]
For every nonzero \(a\in H^1(C,\mathbf F_2)\), the values
\(q_\kappa(a)\), as \(q_\kappa\) ranges over the six odd quadratic
forms, are one four times and zero twice.  This follows by the
transitivity of \(\operatorname{Sp}_4(\mathbf F_2)\) on nonzero vectors
and a direct count for one symplectic basis vector.  Hence the left side
of \eqref{eq:six-parities} can be constant only if \(w_1(W)=0\), and
then \(w_2(W)=1\).

Finally, the determinant of the augmentation representation is the sign
of the five-point permutation action.  Thus \(w_1(W)=0\) forces
\(G\subset A_5\), and \eqref{eq:degree-five-groups} gives \(G=A_5\).
\end{proof}

\begin{proposition}[The remaining theta obstruction]
\label{prop:A5-theta-obstruction}
Let \(\rho\colon\pi_1(C)\twoheadrightarrow A_5\) be a genus-two
surface representation whose natural oriented augmentation bundle
\(W_\rho\) has \(w_2=1\).  Let \(\Theta\simeq C\subset J^1(C)\) be the
Abel theta divisor, and let \(s_\rho\) be the determinant-of-cohomology
theta section of \(W_\rho\), allowing \(s_\rho=0\).  After choosing the
standard theta-line identification up to scalar and a nonzero Wronskian
normalization, there is a unique one-form
\(\omega_\rho\in H^0(C,K_C)\) such that
\[
 s_\rho|_\Theta=\operatorname{Wr}_C\,\omega_\rho,
\tag{6.23}\label{eq:theta-obstruction}
\]
where a nonzero Wronskian
\(\operatorname{Wr}_C\in H^0(C,K_C^3)\) has been chosen; it is
well-defined up to scalar and has zero divisor equal to the six
Weierstrass points.  The form \(\omega_\rho\) is unique relative to this
normalization, and its vanishing is intrinsic.  The
associated degree-five cover is Prill-exceptional if and only if
\(\omega_\rho=0\).

Moreover, it suffices to verify the Prill inequality at three distinct
non-Weierstrass points: three such zeros force
\(\omega_\rho=0\), since a nonzero holomorphic one-form on a genus-two
curve has a zero divisor of degree two.
\end{proposition}

\begin{proof}
For an oriented orthogonal bundle of rank four, Beauville's parity
formula and theta-map lemma \cite[(1.3), Lemma 1.4]{BeauvilleOrthogonal}
place the theta section in the odd eigenspace
\(H^0(J^1(C),\cO(4\Theta))^-\) when \(w_2=1\).  In particular,
\(h^0(C,W_\rho(w_i))\) is odd at every Weierstrass point, so
\(s_\rho|_\Theta\) vanishes at all six points.  Adjunction gives
\(\cO_{J^1(C)}(\Theta)|_\Theta\simeq K_C\), hence
\(s_\rho|_\Theta\in H^0(C,K_C^4)\).  After choosing a nonzero
Wronskian normalization, division gives the corresponding unique
\(\omega_\rho\in H^0(C,K_C)\) in
\eqref{eq:theta-obstruction}.

By the defining property of the determinant-of-cohomology section,
\[
 s_\rho(\cO_C(p))=0
 \quad\Longleftrightarrow\quad
 H^0(C,W_\rho(p))\neq0.
\]
Thus every fibre moves exactly when the
restriction in \eqref{eq:theta-obstruction} is identically zero, which
is equivalent to \(\omega_\rho=0\).  The final assertion follows from
\(\deg K_C=2\).
\end{proof}

\begin{lemma}[Equivariant form of the obstruction]
\label{lem:equivariant-obstruction}
Let a finite group \(H\) act on a genus-two curve \(C\), and suppose
that the oriented orthogonal bundle \(W_\rho\) in
\cref{prop:A5-theta-obstruction} is \(H\)-equivariant.  Put
\[
 V_C=H^0(C,K_C),\qquad
 D_\rho=\det R\Gamma(C,W_\rho),
\]
and let \(\Delta_\rho\) be the one-dimensional \(H\)-character carried
by the \(H\)-linearized trivial determinant line
\(\det W_\rho\simeq\cO_C\).  Define the character line
\[
 \Xi_\rho=
 D_\rho^{-1}\otimes\Delta_\rho^{-1}\otimes\det V_C.
\tag{6.24}\label{eq:Xi-character}
\]
Then the intrinsic obstruction is an invariant tensor
\[
 \Omega_\rho\in(V_C\otimes\Xi_\rho)^H.
\tag{6.25}\label{eq:equivariant-obstruction-space}
\]
In particular, if the invariant space in
\eqref{eq:equivariant-obstruction-space} is zero, then the associated
degree-five cover is Prill-exceptional.
\end{lemma}

\begin{proof}
Let \(a\colon C\to J^1(C)\) be \(p\mapsto\cO_C(p)\), and use the
inverse determinant-of-cohomology convention for the theta line
\(\lambda_\rho\).  The exact sequence
\[
 0\longrightarrow W_\rho\longrightarrow W_\rho(p)
 \longrightarrow W_{\rho,p}\otimes T_pC\longrightarrow0
\]
gives the \(H\)-equivariant identity
\[
 a^*\lambda_\rho\simeq
 D_\rho^{-1}\otimes\Delta_\rho^{-1}\otimes K_C^4.
\tag{6.26}\label{eq:abel-theta-line-equivariant}
\]
The basis-free canonical Wronskian is a tensor
\[
 \mathfrak W_C\in
 H^0(C,K_C^3)\otimes(\det V_C)^{-1}.
\]
Dividing the invariant restriction \(a^*s_\rho\) by
\(\mathfrak W_C\) therefore gives precisely an invariant element of
\(V_C\otimes\Xi_\rho\).  After bases are chosen in the character
lines, this is the one-form \(\omega_\rho\) of
\cref{prop:A5-theta-obstruction}; only its vanishing is independent of
those choices.
\end{proof}

\begin{theorem}[A degree-five Prill cover and the genus-two minimum]
\label{thm:C5-degree-five}
Let
\[
 C:\ y^2=x^5-1.
\]
There exists a connected degree-five \'etale cover
\(f\colon X\to C\) such that
\[
 h^0\!\left(X,\cO_X(f^{-1}(p))\right)\geq2
 \qquad\text{for every }p\in C.
\tag{6.27}\label{eq:C5-Prill-condition}
\]
Its monodromy group is \(A_5\) in the natural five-point action,
\(g(X)=6\), and its \(A_5\)-Galois closure has degree \(60\) and genus
\(61\).  Consequently
\[
 \min\!\left\{\deg f'\ \middle|\
 \begin{array}{c}
 C'\text{ is a genus-two curve},\\
 f'\colon X'\to C'\text{ is connected and Prill-exceptional}
 \end{array}
 \right\}=\boxed{5}.
\tag{6.28}\label{eq:genus-two-global-minimum}
\]
\end{theorem}

\begin{proof}
Fix a primitive fifth root \(\zeta\), and let
\(\sigma(x,y)=(\zeta x,y)\).  The fixed points are
\(P_\pm=(0,\pm i)\) and the point \(P_\infty\) at infinity.  The
tangent multipliers of \(\sigma\) are \(\zeta,\zeta,\zeta^2\): at
\(P_\pm\) use \(x\), while at infinity one may use
\(t=x^2/y\), for which \(\sigma(t)=\zeta^2t\).  Hence
\(C/\langle\sigma\rangle\simeq\mathbf P^1\) has signature
\((0;5,5,5)\), with orbifold epimorphism
\[
 \Delta(5,5,5)=
 \langle x_1,x_2,x_3\mid x_i^5=x_1x_2x_3=1\rangle
 \longrightarrow\langle\sigma\rangle,
 \qquad (x_1,x_2,x_3)\longmapsto(\sigma,\sigma,\sigma^3).
\tag{6.29}\label{eq:C5-orbifold-map}
\]

In \(SL_2(\mathbf F_5)\), take
\[
 U_1=\begin{pmatrix}1&1\\0&1\end{pmatrix},\qquad
 U_2=\begin{pmatrix}3&1\\1&4\end{pmatrix},\qquad
 U_3=\begin{pmatrix}1&0\\1&1\end{pmatrix}.
\tag{6.30}\label{eq:C5-SL-triple}
\]
Direct multiplication gives
\[
 U_1^5=U_2^5=U_3^5=I,\qquad U_1U_2U_3=-I.
\tag{6.31}\label{eq:C5-SL-relations}
\]
Moreover \(U_1,U_3\) generate \(SL_2(\mathbf F_5)\), since they are
the elementary upper and lower unipotent matrices.  If
\(\bar U_i\) denotes the image in
\(PSL_2(\mathbf F_5)\simeq A_5\), then
\[
 x_1\longmapsto(\bar U_1,\sigma),\qquad
 x_2\longmapsto(\bar U_2,\sigma),\qquad
 x_3\longmapsto(\bar U_3,\sigma^3)
\tag{6.32}\label{eq:C5-extension-map}
\]
defines a homomorphism onto \(A_5\times C_5\).  Indeed, both
projections are surjective, and \(A_5\) and \(C_5\) have no common
nontrivial quotient.  Its restriction to the index-five surface
kernel in \eqref{eq:C5-orbifold-map} is therefore surjective onto
\(A_5\).  The natural five-point action gives a connected degree-five
\'etale cover and an equivariant augmentation bundle
\(W_\rho\).

I next compute its spin obstruction.  The pullback of
\(\operatorname{Spin}(4)\to SO(4)\) under the natural augmentation
embedding is the Schur cover
\[
 2\!\cdot\!A_5\simeq SL_2(\mathbf F_5).
\]
It is non-split: a splitting would embed the simple group \(A_5\) in
\(\operatorname{Spin}(4)=SU(2)\times SU(2)\), so one projection would
embed \(A_5\) in \(SU(2)\), which is impossible because \(SU(2)\) has
only one element of order two whereas \(A_5\) has fifteen.
Each order-five element of \(PSL_2(\mathbf F_5)\) has a unique lift
whose fifth power is \(I\); the other lift has fifth power \(-I\).
Thus the orbifold relations force the three lifts in
\eqref{eq:C5-SL-triple}, whose product is \(-I\) by
\eqref{eq:C5-SL-relations}.  The orbifold representation has a nonzero
spin-lifting class.  Restriction to the surface subgroup preserves
nonvanishing: with \(\mathbf F_2\)-coefficients,
\(\operatorname{cor}\circ\operatorname{res}\) is multiplication by
the index \(5\), hence is the identity.  Consequently
\[
 \langle w_2(W_\rho),[C]\rangle=1.
\tag{6.33}\label{eq:C5-w2-one}
\]

It remains to kill the holomorphic obstruction.  Write
\(\chi(\sigma)=\zeta\).  Pullback acts on canonical forms by
\[
 V_C=H^0(C,K_C)
 =\left\langle\frac{dx}{y},\frac{x\,dx}{y}\right\rangle
 \simeq\chi\oplus\chi^2,\qquad \det V_C=\chi^3.
\tag{6.34}\label{eq:C5-Hodge-character}
\]
The augmentation character has value \(-1\) on every nontrivial
five-cycle.  At infinity the action of \(\sigma\) on the fibre is the
square of the third inertia element, which is still a nontrivial
five-cycle.  Since \(H^0(C,W_\rho)=0\), Serre duality gives
\(H^1(C,K_C\otimes W_\rho)=0\).  The holomorphic Lefschetz formula
\cite{AtiyahBottLefschetz},
with pullback linearization and denominator \(1-\lambda\), gives for
\(1\leq k\leq4\)
\[
 \operatorname{Tr}\!\left(\sigma^k\mid
 H^0(C,K_C\otimes W_\rho)\right)
 =-\left(
 \frac{2\zeta^k}{1-\zeta^k}
 +\frac{\zeta^{2k}}{1-\zeta^{2k}}
 \right).
\tag{6.35}\label{eq:C5-Lefschetz-trace}
\]
For a nontrivial fifth root \(q\),
\[
 \frac{2q}{1-q}+\frac{q^2}{1-q^2}=q+q^2-1.
\]
It follows that
\[
 H^0(C,K_C\otimes W_\rho)
 \simeq\mathbf1^{\oplus2}\oplus\chi^3\oplus\chi^4.
\tag{6.36}\label{eq:C5-KW-character}
\]
With the convention
\(D_\rho=\det H^0(C,W_\rho)\otimes
\det H^1(C,W_\rho)^{-1}\), Serre duality and
\eqref{eq:C5-KW-character} give
\[
 D_\rho=\chi^2.
\]
The representation \(A_5\times C_5\to SO(4)\) factors through
\(A_5\), so its determinant character is trivial and
\(\Delta_\rho=\mathbf1\).
Therefore
\[
 \Xi_\rho=D_\rho^{-1}\otimes\Delta_\rho^{-1}\otimes\det V_C
 =\chi,
 \qquad
 V_C\otimes\Xi_\rho=\chi^2\oplus\chi^3.
\tag{6.37}\label{eq:C5-Xi-character}
\]
The latter representation has no invariants.  By
\cref{lem:equivariant-obstruction}, \(\Omega_\rho=0\), and then
\cref{prop:A5-theta-obstruction} gives
\(H^0(C,W_\rho(p))\neq0\) for every \(p\in C\).  Since
\[
 f_*\cO_X\simeq\cO_C\oplus W_\rho,
\]
this is exactly \eqref{eq:C5-Prill-condition}.

The genus statements follow from the unramified Riemann--Hurwitz
formula.  Finally, \cref{cor:cover-lower-bound} gives degree at least
\(2+3=5\) for every connected genus-two Prill-exceptional cover, so
the constructed degree-five cover proves
\eqref{eq:genus-two-global-minimum}.
\end{proof}

\begin{remark}[Reproducible finite certificates]
The supplementary GAP certificate and independent Python character
check verify the finite group generation, the lifted surface relation,
and the \(C_5\)-character decomposition used above.  The construction only
forces the generalized theta divisor to contain the Abel curve
\(\Theta\simeq C\); it does not assert vanishing over all of
\(\operatorname{Pic}^1(C)\).
\end{remark}

\section{Consequences for generic generation and horizontal vanishing}

I record the exact implications for the questions formulated in
\cite[\S6.3]{LittMotives}.  Throughout this section \(g\geq3\), and I
take \(n=g\) in \cref{thm:uniform}.  Write \(V\) for the resulting
rank-\(g^g\) flat bundle.

\begin{proposition}[Failure of generic global generation]\label{prop:ggg-failure}
The fixed irreducible finite-unitary representation dual to
\(\rho_{g,g}\) has the following property for every marked complex
structure:
\[
 K_C\otimes V^\vee
\quad\text{is not generated by global sections at any point of }C.
\tag{7.1}\label{eq:nowhere-ggg}
\]
Consequently, Conjectures 6.3.5 and 6.3.11 of
\cite{LittMotives} are false.
\end{proposition}

\begin{proof}
Since \(H^0(V)=0\), Serre duality gives
\(H^1(K_C\otimes V^\vee)=0\).  From
\[
 0\longrightarrow K_C\otimes V^\vee(-p)
 \longrightarrow K_C\otimes V^\vee
 \longrightarrow (K_C\otimes V^\vee)|_p
 \longrightarrow0,
\]
the dual of the evaluation cokernel at \(p\) is
\[
 H^1(K_C\otimes V^\vee(-p))^\vee
 \simeq H^0(V(p)).
\tag{7.2}\label{eq:evaluation-dual}
\]
This space is nonzero for every \(p\) by \cref{thm:uniform}.  Hence
evaluation fails to be surjective at every point, for every complex
structure.  The cited conjectures predict generic, respectively global,
generation for every fixed unitary representation on a generic complex
structure, so this fixed representation refutes both predictions.
\end{proof}

\begin{proposition}[Failure of horizontal generic vanishing]\label{prop:HGV-failure}
In the unpunctured case \(D=\varnothing\), the fixed monodromy
\(\rho_{g,g}\) satisfies
\[
 H^0(C,V\otimes L)\neq0
 \quad\text{for every marked complex structure and every }
 L\in\Pic^1(C).
\tag{7.3}\label{eq:horizontal-failure}
\]
Thus both the weak and strong forms of Conjecture 6.3.17 in
\cite{LittMotives} fail.
\end{proposition}

\begin{proof}
Equation \eqref{eq:horizontal-failure} is
\cref{thm:uniform} with \(n=g\).  Apply Conjecture 6.3.17 with
\(D=\varnothing\).  Its nondecreasing function has value \(1\) at
\(g=3\), hence value at least \(1\) for every \(g\geq3\).  Taking the
degree-one divisor \(Z=p\), the nonvanishing
\(H^0(C,V(p))\neq0\) for every \(p\), and for every isomonodromic
deformation, contradicts both the weak and strong forms.
\end{proof}

\begin{corollary}[Status of the adjacent questions]\label{cor:question-status}
The construction has the following consequences.
\begin{enumerate}[label=\textup{(\roman*)}]
\item It answers Prill's problem affirmatively for every base genus
\(g\geq2\), with every fibre moving.
\item It disproves the unpunctured specialization, Conjecture 6.3.5, of
the parabolic generic-generation Conjecture 6.3.4.
\end{enumerate}
\end{corollary}

\begin{proof}
The first statement is \cref{cor:prill-all-genus}.  The second is
\cref{prop:ggg-failure} together with the logical relation among the
conjectures stated in \cite[\S6.3]{LittMotives}.
\end{proof}

\begin{remark}[Scope of the consequences]\label{rem:adjacent-scope}
The construction is unpunctured and supplies no conclusion for the
genuinely punctured cases of Conjecture 6.3.4.  No implication is
established here from these support-locus statements to a
Putman--Wieland or Ivanov counterexample, a constant Prym variation, or
an isotrivial Jacobian factor; see \cref{rem:no-PW-converse}.
\end{remark}

\section{Abelian support blocks and Lagrangian theta descent}

This section isolates the construction from the special choice
\(A=J(C)\).  It produces abelian blocks in support loci whenever a curve
maps homologically onto an abelian variety.

Let \(A\) be a complex abelian variety of dimension \(b\), let
\(\widehat A=\Pic^0(A)\), and let \(N\) be an ample line bundle on \(A\).
For \(a\in A=\Pic^0(\widehat A)\), write
\(P_a=\mathcal P|_{\widehat A\times\{a\}}\) for the corresponding
topologically trivial line bundle on \(\widehat A\).
Write
\[
 \phi_N\colon A\longrightarrow\widehat A,\qquad
 r=h^0(A,N),\qquad K(N)=\ker\phi_N.
\tag{8.1}\label{eq:general-polarization}
\]
Thus \(|K(N)|=r^2\).  Let \(N_0\) be the exponent of \(K(N)\), and define
the Fourier--Mukai bundle
\[
 F_N=\Phi_{\mathcal P}(N^{-1})[b]
\quad\text{on }\widehat A.
\tag{8.2}\label{eq:general-FM}
\]
For a morphism \(u\colon C\to\widehat A\), put \(R_{u,N}=u^*F_N\).

\begin{theorem}[Lagrangian theta descent and abelian support]\label{thm:abelian-block}
Assume that
\[
 u_*\colon H_1(C,\bZ)\twoheadrightarrow H_1(\widehat A,\bZ)
\tag{8.3}\label{eq:H1-surjective}
\]
is surjective.  Then the following statements hold.
\begin{enumerate}[label=\textup{(\roman*)}]
\item The bundle \(R_{u,N}\) is stable of rank \(r\).  Its projective
monodromy is the irreducible theta representation of the finite
symplectic module \(K(N)\), with the central weight fixed by the dual
section factor in \eqref{eq:general-mukai-splitting}.
\item Let \(\beta_N\in H^2(\widehat A,\bQ)\) be defined by
\[
 \phi_N^*\beta_N=c_1(N).
\tag{8.4}\label{eq:beta-definition}
\]
Then
\[
 \mu(R_{u,N})=\int_Cu^*\beta_N.
\tag{8.5}\label{eq:general-slope}
\]
\item If
\[
 B_u=\im\!\left(
 u^*\colon A=\Pic^0(\widehat A)\longrightarrow\Pic^0(C)\right),
\tag{8.6}\label{eq:Bu}
\]
then
\[
 B_u\subset V_0^0(R_{u,N}).
\tag{8.7}\label{eq:abelian-support}
\]
Here
\(V_0^0(R_{u,N})=\{\alpha\in\Pic^0(C):
H^0(C,R_{u,N}\otimes\alpha)\neq0\}\).
\item After choosing a linearized maximal isotropic subgroup of
\(K(N)\), there is a connected degree-\(r\) \'etale cover
\(q\colon C_L\to C\) and a line bundle \(\mathcal L\) on \(C_L\) such
that
\[
 R_{u,N}\simeq q_*\mathcal L,\qquad
 H^i(C,R_{u,N}\otimes A_C)
 \simeq H^i(C_L,\mathcal L\otimes q^*A_C)
\tag{8.8}\label{eq:general-line-descent}
\]
for every line bundle \(A_C\) on \(C\) and \(i=0,1\).  Moreover,
\[
 \deg\mathcal L=\deg R_{u,N}
 =r\,\mu(R_{u,N}).
\tag{8.9}\label{eq:line-degree}
\]
\end{enumerate}
\end{theorem}

\begin{proof}
Mukai's isogeny formula, with its theta-group equivariance, gives
\[
 \phi_N^*F_N\simeq
 (-1_A)^*N\otimes H^0(A,N)^\vee.
\tag{8.10}\label{eq:general-mukai-splitting}
\]
The fibre product
\[
 \widetilde C=C\times_{\widehat A,\phi_N}A
\]
is connected by \eqref{eq:H1-surjective}.  Pulling
\eqref{eq:general-mukai-splitting} to \(\widetilde C\) shows that
\(R_{u,N}\) has the stated projective monodromy.

The same identity makes the pullback of \(R_{u,N}\) a direct sum of
copies of one line bundle, so \(R_{u,N}\) is semistable.  To prove
stability, suppose a proper subbundle has the same slope.
After pullback to \(\widetilde C\) and tensoring by the inverse of the
line factor in \eqref{eq:general-mukai-splitting}, it becomes a
degree-zero subbundle \(G\) of
\(\cO_{\widetilde C}\otimes H^0(A,N)^\vee\).  Such a subbundle is
constant.  Indeed, \(G^\vee\) is a globally generated degree-zero
quotient of a trivial bundle.  Choosing \(\rk G\) generically independent
sections produces a torsion quotient of length \(\deg G^\vee=0\), so
\(G^\vee\simeq\cO_{\widetilde C}^{\oplus\rk G}\), and the inclusion of
\(G\) is induced by a constant subspace
\(W\subset H^0(A,N)^\vee\).  Equivariance of the deck descent forces
\(W\) to be
invariant under the full projective Schr\"odinger action, contrary to
irreducibility unless \(W=0\) or \(W=H^0(A,N)^\vee\).  This proves (i).

Taking first Chern classes in
\eqref{eq:general-mukai-splitting} gives
\[
 \phi_N^*\!\left(\frac{c_1(F_N)}r\right)=c_1(N).
\]
Pullback to \(C\) proves (ii).

For (iii), Mukai inversion \cite{Mukai} gives
\[
 \widehat\Phi_{\mathcal P}(F_N)\simeq(-1_A)^*N^{-1}
\]
in degree zero.  Thus \(F_N\) satisfies \(\mathrm{IT}_0\), has Euler
characteristic one, and base change gives
\[
 h^0(\widehat A,F_N\otimes P_a)=1,\qquad
 H^i(\widehat A,F_N\otimes P_a)=0\quad(i>0)
\]
for every \(a\in A\).  These one-dimensional spaces are the fibres of
the transform line.  After trivializing that line over a nonempty open
subset of \(A\), base change produces an evaluation morphism.  It is
nonzero somewhere, since each fibre is a nonzero global section; hence it
is nonzero at some point \(x_0\in\widehat A\).  Since \(F_N\) is simple
semihomogeneous, translation by every
\(t\in\widehat A\) has the form
\[
 t_t^*F_N\simeq F_N\otimes P_{\psi(t)}
\tag{8.11}\label{eq:semihomogeneous-translation}
\]
for a suitable \(\psi(t)\in A\).  Translating the preceding family shows
that, for every prescribed \(x\in\widehat A\), there is a nonempty open
set \(U_x\subset A\) on which the corresponding section is nonzero at
\(x\).

Fix \(p\in C\) and take \(x=u(p)\).  For \(a\in U_x\), restriction gives
\[
 H^0(C,R_{u,N}\otimes u^*P_a)\neq0.
\]
The locus of \(a\) satisfying this inequality is closed by upper
semicontinuity.  It contains a nonempty open subset of the irreducible
variety \(A\), and hence it is all of \(A\).  This proves
\eqref{eq:abelian-support}.  This continuous-generation argument is the
general ample-polarization form of Raynaud's construction
\cite[\S3]{Raynaud}.

Finally, choose a maximal isotropic theta subgroup.  Standard
Lagrangian descent for theta groups, in the form used in
\cite[Proposition 2.1(2)]{Schneider}, gives an abelian variety
\(\widehat A_L\), an isogeny
\(\tau\colon\widehat A_L\to\widehat A\) of degree \(r\), and a line
bundle \(N_L\) such that
\[
 F_N\simeq\tau_*N_L.
\]
The base change \(C_L=C\times_{\widehat A}\widehat A_L\) is connected by
\eqref{eq:H1-surjective}.  Flat base change proves
\eqref{eq:general-line-descent}.  Since \(q\) is \'etale,
Riemann--Roch gives \(\deg\mathcal L=\deg q_*\mathcal L\), proving
\eqref{eq:line-degree}.
\end{proof}

\begin{theorem}[Integral-slope Heisenberg untwisting]\label{thm:general-untwisting}
In the setting of \cref{thm:abelian-block}, suppose
\[
 \mu=\mu(R_{u,N})\in\bZ.
\tag{8.12}\label{eq:integral-slope}
\]
Assume in addition that \(K(N)\neq0\), equivalently \(N_0\geq2\).
Then there is a finite irreducible unitary local system \(V\) of rank
\(r\), a line bundle \(M\in\Pic^\mu(C)\), and a connected \'etale cover
\[
 f\colon Z\longrightarrow C,\qquad \deg f=N_0r,
\tag{8.13}\label{eq:general-cover-degree}
\]
such that
\[
 R_{u,N}\simeq V\otimes M,\qquad
 f_*\cO_Z\simeq\cO_C\oplus V\oplus E,
\tag{8.14}\label{eq:general-splitting}
\]
Writing
\(M+B_u:=\{M\otimes\alpha:\alpha\in B_u\}\subset\Pic^\mu(C)\), one has
\[
 h^0(Z,f^*L)\geq h^0(C,L)+1
\quad\text{for every }L\in M+B_u.
\tag{8.15}\label{eq:block-pullback-sections}
\]
The Galois closure has degree \(N_0r^2\).  Equivalently,
\[
 f^*(M+B_u)\subset W_{N_0r\mu}^0(Z).
\tag{8.16}\label{eq:rank-one-block}
\]
\end{theorem}

\begin{proof}
For a projectively flat bundle \(E\) of rank \(r\) and degree \(d\), the
exact sequence
\[
 1\longrightarrow\mu_r\longrightarrow SL_r
 \longrightarrow PGL_r\longrightarrow1
\]
associates to \(\bP(E)\) an obstruction
\(o(E)\in H^2(C,\mu_r)\).  Under
\(H^2(C,\mu_r)\simeq\bZ/r\), this class is \(d\bmod r\): twisting a
\(GL_r\)-lift by a line bundle changes its determinant degree by a
multiple of \(r\), and determinant reduction gives the converse.  Under
the embedding \(\mu_r\subset\bC^*\), evaluation on \([C]\) is therefore
\[
 \exp(\pm2\pi i\,d/r),
\tag{8.17}\label{eq:degree-obstruction}
\]
where the sign depends only on the commutator convention.  Here
\(d/r=\mu\), so \eqref{eq:integral-slope} kills the surface obstruction.
The projective theta class is determined by its commutator because
\(H^2(K(N),\bC^*)\simeq
\Hom(\wedge^2K(N),\bC^*)\).  It may therefore be represented by the
split finite model \(H_{\mathrm{sp}}(K(N))\) constructed in
\cref{thm:finite-symplectic-untwisting}:
\[
 1\longrightarrow\mu_{N_0}\longrightarrow H_{\mathrm{sp}}(K(N))
 \longrightarrow K(N)\longrightarrow1.
\]
Choosing lifts of surface generators in this model, the vanishing
obstruction makes their product of commutators one; hence
\cref{lem:surface-lift} gives a homomorphism
\(\pi_1(C)\to H_{\mathrm{sp}}(K(N))\).  Its Schr\"odinger representation, with the
central weight matching \eqref{eq:general-mukai-splitting}, defines
\(V\).

The finite lift is surjective.  Indeed, its projective image is all of
\(K(N)\), while commutators of lifts generate \(\mu_{N_0}\) because the
theta pairing is perfect and has exponent \(N_0\).

Choose \(M_0\in\Pic^\mu(C)\).  The bundles \(R_{u,N}\) and
\(V\otimes M_0\) have the same projectivization and the same degree.
By \cref{lem:scalar-transport},
\(R_{u,N}\simeq V\otimes M_0\otimes\alpha\) for some
\(\alpha\in\Pic^0(C)\).  Setting \(M=M_0\otimes\alpha\) gives the first
identity in \eqref{eq:general-splitting}.

Equation \eqref{eq:abelian-support} now says
\[
 H^0(C,V\otimes L)\neq0
\quad(L\in M+B_u).
\tag{8.18}\label{eq:translated-support}
\]
Apply \cref{thm:finite-symplectic-untwisting} to this surjection.  It
produces the degree-\(N_0r\) cover, places
\(\cO_C\oplus V\) in its permutation local system, and gives the
core-free Galois closure of degree \(N_0r^2\).  The projection formula
and \eqref{eq:translated-support} prove
\eqref{eq:block-pullback-sections}; its rank-one reformulation is
\eqref{eq:rank-one-block}.
\end{proof}

\begin{corollary}[Prym-direction Brill--Noether block]\label{cor:prym-block}
Let \(P\subset J(Y)\) be an abelian subvariety of dimension \(b\),
endowed with a principal polarization \(\Xi\) not assumed to be the
restriction of the theta polarization of \(J(Y)\).  Use the Jacobian
principal polarization to identify \(J(Y)\simeq J(Y)^\vee\), and let
\[
 q_P\colon J(Y)\longrightarrow P^\vee
\]
be dual to the inclusion, and let
\(u_P=q_P\circ\operatorname{AJ}_Y\).  If \(\widehat\Xi\) is the dual
principal class, assume
\[
 D=\deg(u_P^*\widehat\Xi)\geq2.
\tag{8.19}\label{eq:Prym-D}
\]
Then there is an \(M\in\Pic^1(Y)\) and a connected \'etale cover
\[
 f\colon Z\longrightarrow Y,\qquad
 \deg f=D^{b+1},
\tag{8.20}\label{eq:Prym-cover-degree}
\]
such that
\[
 f^*(M+u_P^*P)\subset W_{D^{b+1}}^0(Z).
\tag{8.21}\label{eq:Prym-block}
\]
Its Heisenberg Galois closure has degree \(D^{2b+1}\).
\end{corollary}

\begin{proof}
The kernel of \(q_P\) is connected, so \(u_P\) is surjective on integral
first homology.  Apply \cref{thm:abelian-block} with
\(A=P\) and \(N=\Xi^{\otimes D}\).  Then
\[
 r=D^b,\qquad K(N)=P[D],\qquad N_0=D.
\]
Since
\[
 \beta_N=\frac1D\,c_1(\widehat\Xi),
\]
\cref{eq:general-slope,eq:Prym-D} give
\(\mu(R_{u_P,N})=1\).  Now apply
\cref{thm:general-untwisting}.
\end{proof}

\begin{remark}[Prym boundary]\label{rem:Prym-boundary}
If \(P\) is a proper abelian subvariety of \(J(Y)\), the translate
\(M+u_P^*P\) is only a proper abelian block in \(\Pic^1(Y)\).  It need
not contain the Abel curve \(p\mapsto\cO_Y(p)\).  Thus
\cref{cor:prym-block} is a Brill--Noether statement, not a full Prill
statement and not a claim about finite higher-Prym monodromy.
The statement applies in particular to principally polarized Prym
varieties; no principal polarization is asserted for an arbitrary
ramified Prym.
\end{remark}

\section{Exact boundaries of the construction}

I finish by isolating the hypotheses that cannot be removed formally.

\begin{proposition}[Integrality is the flat lifting condition]\label{prop:integrality-boundary}
Let \(E\) be a projectively flat bundle of rank \(r\) and degree \(d\)
whose projective monodromy is a finite theta representation.  Its
projective representation has a flat linear lift on \(C\) if and only if
\[
 r\mid d.
\tag{9.1}\label{eq:integrality-iff}
\]
If the scalar obstruction has order \(s\), then pullback along a
connected finite \'etale cover of degree \(\ell\) kills it if and only
if \(s\mid\ell\).
\end{proposition}

\begin{proof}
The obstruction class is \(d\bmod r\in H^2(C,\mu_r)\), as proved above,
so it vanishes exactly when \(r\mid d\).  If \(p\colon C'\to C\) has
degree \(\ell\), then
\(p_*[C']=\ell[C]\); hence evaluation of \(p^*o(E)\) is the
\(\ell\)-th power of the original scalar.  A scalar of exact order
\(s\) is killed precisely when \(s\mid\ell\).
\end{proof}

This proposition concerns a flat lift of the projective representation;
the bundle \(E\) itself is flat only when \(d=0\).  A lift inside the
finite split theta model is obtained, when needed, by the construction
in \cref{thm:finite-symplectic-untwisting}.

For the Raynaud projective bundle at level \(n\), the scalar is
\(\zeta^{-g}\), of order
\[
 \frac n{\gcd(n,g)}.
\tag{9.2}\label{eq:raynaud-obstruction-order}
\]
Thus \(n\mid g\) is not a convenient sufficient assumption; it is the
exact condition for a flat scalar lift on the original surface.  Killing
the scalar after a base change does not transfer the full support
property from \(\Pic(C)\) to the full Picard variety of the new curve.

\begin{remark}[Marked versus characteristic]\label{rem:marked-characteristic}
The homomorphism \(\varphi_{g,n}\) and every cover in the tower are fixed
after choosing a marking.  Their kernels are not claimed to be
characteristic subgroups of \(\pi_g\).  The projective theta class is
symplectic, but linear lifts can be permuted by scalar characters under
the mapping class group.
\end{remark}

\begin{remark}[Even and composite levels]\label{rem:even-composite}
No primeness or parity hypothesis on \(n\) appears in the construction.
The trace calculation, isotropic order bound, and exact
degree--multiplicity law are valid over \(\bZ/n\).  Even-level quadratic
refinements become relevant only when extending the symplectic group to
a prescribed metaplectic action.  The fixed marked projective class is
already determined by its commutator pairing.
\end{remark}

\begin{remark}[The three geometric outputs]\label{rem:three-outputs}
The method separates three different kinds of information:
\[
 \begin{array}{ccl}
 \text{theta-equivariant Fourier--Mukai splitting}
   &\Longrightarrow&\text{finite projective monodromy},\\
 \text{integral slope and scalar transport}
   &\Longrightarrow&\text{a fixed flat support block},\\
 \text{Lagrangian induction and cyclic untwisting}
   &\Longrightarrow&\text{rank-one geometry on a finite cover}.
 \end{array}
\tag{9.3}\label{eq:three-outputs}
\]
The last step is what lowers the ordinary cover degree from
\(n^{2g+1}\) to \(n^{g+1}\).  The middle step is what makes the
construction uniform over marked complex structures.  Neither step
alone implies fixed-part monodromy.
\end{remark}

\begin{thebibliography}{99}
\small

\bibitem{ACGH}
E.~Arbarello, M.~Cornalba, P.~A.~Griffiths, and J.~Harris,
\emph{Geometry of algebraic curves. Vol. I},
Grundlehren Math. Wiss. \textbf{267}, Springer, New York, 1985.

\bibitem{AtiyahSpin}
M.~F.~Atiyah,
Riemann surfaces and spin structures,
\emph{Ann. Sci. \'Ecole Norm. Sup.} (4) \textbf{4} (1971), 47--62.

\bibitem{AtiyahBottLefschetz}
M.~F.~Atiyah and R.~Bott,
A Lefschetz fixed point formula for elliptic complexes. II.
Applications,
\emph{Ann. of Math.} (2) \textbf{88} (1968), 451--491.

\bibitem{BeauvilleOrthogonal}
A.~Beauville,
Orthogonal bundles on curves and theta functions,
\emph{Ann. Inst. Fourier (Grenoble)} \textbf{56} (2006), 1405--1418.

\bibitem{BrambilaPazGN}
L.~Brambila-Paz, I.~Grzegorczyk, and P.~E.~Newstead,
Geography of Brill--Noether loci for small slopes,
\emph{J. Algebraic Geom.} \textbf{6} (1997), 645--669.

\bibitem{LandesmanLittApplications}
A.~Landesman and D.~Litt,
Applications of the algebraic geometry of the Putman--Wieland
conjecture,
\emph{Proc. Lond. Math. Soc.} (3) \textbf{127} (2023), 116--133.

\bibitem{LandesmanLittPrill}
A.~Landesman and D.~Litt,
Prill's problem,
\emph{Algebraic Geometry} \textbf{11} (2024), 290--295.

\bibitem{LandesmanLittSawin}
A.~Landesman, D.~Litt, and W.~Sawin,
Big monodromy for higher Prym representations,
\emph{Geom. Topol.} \textbf{29} (2025), 2733--2782.

\bibitem{LittMotives}
D.~Litt,
Motives, mapping class groups, and monodromy,
arXiv:2409.02234v1, 2024.

\bibitem{Mukai}
S.~Mukai,
Duality between \(D(X)\) and \(D(\widehat X)\) with its application to
Picard sheaves,
\emph{Nagoya Math. J.} \textbf{81} (1981), 153--175.

\bibitem{MumfordTheta}
D.~Mumford,
\emph{Tata lectures on theta. I},
Progress in Mathematics \textbf{28}, Birkh\"auser, Boston, 1983.

\bibitem{NarasimhanSeshadri}
M.~S.~Narasimhan and C.~S.~Seshadri,
Stable and unitary vector bundles on a compact Riemann surface,
\emph{Ann. of Math.} (2) \textbf{82} (1965), 540--567.

\bibitem{PutmanWieland}
A.~Putman and B.~Wieland,
Abelian quotients of subgroups of the mapping class group and higher
Prym representations,
\emph{J. Lond. Math. Soc.} (2) \textbf{88} (2013), 79--96.

\bibitem{Raynaud}
M.~Raynaud,
Sections des fibr\'es vectoriels sur une courbe,
\emph{Bull. Soc. Math. France} \textbf{110} (1982), 103--125.

\bibitem{Schneider}
O.~Schneider,
Sur la dimension de l'ensemble des points base du fibr\'e d\'eterminant
sur \(SU_C(r)\),
\emph{Ann. Inst. Fourier (Grenoble)} \textbf{57} (2007), 481--490.

\bibitem{SerreFiniteGroups}
J.-P.~Serre,
\emph{Linear representations of finite groups},
Graduate Texts in Mathematics \textbf{42}, Springer, New York, 1977.

\end{thebibliography}

\end{document}
